%% 第26回編集委員会（物理第4巻）：組 電W6
%% p4-11.tex  電-11　電力を運ぶ3本の線：交流電力と三相交流（★2〜5・3段）
%% 中身は計画書 tools/phys4-plan.md の B「電-11」。公式の持ち主は A56（交流の電力・力率）・A78（有効・無効・皮相電力、力率改善）・A79（三相交流）、
%%   A17 の「高い電圧で送ると損失が小さい」の式（⑥。はてな?の答えは電-0 ③、変圧器は電-12 ⑤）。
%% 数値はすべて Python で計算した：checks/物理第4巻/電-11/calc_p4_11.py（出力 calc_p4_11_out.txt）。図の線は pgfplots の式で、点は同じ出力で確かめた。
%% 記号（この章）：瞬時の電圧・電流 v・i、実効値 V・I（数学 4-9 の V_e・I_e と同じ量。Y-3）、電源の相の起電力 e_a・e_b・e_c（実効値 E）、
%%   電流のおくれの角 φ（電位はこの章で使わない）、瞬時電力 p、有効電力 P〔W〕・無効電力 Q〔var〕・皮相電力 S〔VA〕（電荷・面積はこの章で使わない）、
%%   力率 cosφ、線間電圧 V・相電圧 V_p、線電流 I・相電流 I_p、送電線の抵抗 R、リアクタンス X、インピーダンス Z、
%%   対称座標の a＝e^{j2π/3}（6-9 の1の3乗根。相の名前の添字 a・b・c は立体）、極の数は記号にしない（瞬時電力 p とぶつけない。極数 p は電-12）。
%% 数学の本が引いている所（変えない数・言い方）：4-9⑧（p＝V_0I_0 sin² の平均、実効値 V_e、電流が60°おくれると半分、力率）、
%%   5-15⑨（100 V の山 141 V）、6-8⑦（共役）、6-9:541・856（三相の和は0、回転磁界「物理シリーズ『電磁気学』」）→ ⑦⑨、
%%   6-10:578（三相のフェーザ）→ ⑦。
%% 前の巻との約束：第1巻の巻末 ph1-back:57「電力（電-11）」は電-4 に直す（計画書 P10。事務局へ）。この章では電力 P のカードを置かない。
%% 出典（本文で使った事実。手引き BRIEF-phys4 2章の16 に従い、書くときに確かめた）
%%   ・家庭の電圧：電気事業法施行規則 第38条（標準電圧 100 V は「百一ボルトの上下六ボルトを超えない値」、200 V は「二百二ボルトの上下二十ボルト」）。
%%       e-Gov 法令 API https://laws.e-gov.go.jp/api/1/lawdata/407M50000400077 で本文を読んだ（11月5日）。
%%   ・磁束密度 200 μT：電気設備に関する技術基準を定める省令 第27条の2（変電所・電線路などの付近の、人が占める空間の磁束密度の平均値、商用周波数）。
%%       https://laws.e-gov.go.jp/api/1/lawdata/409M50000400052 で本文を読んだ。
%%   ・WHO 環境保健クライテリア 238（2007）超低周波電磁界、環境省版の日本語訳 https://www.env.go.jp/chemi/electric/material/ehc238_j.pdf
%%       （1.1 要約の「慢性影響」：0.3〜0.4 μT 以上の長いばく露と小児白血病の疫学研究、「因果関係があると考えるほどには証拠は強くないが、関心を残すには十分に強い」、
%%       実験室の証拠はしくみを支持できていない。1.1.4：家庭の平均 0.05〜0.1 μT 程度、電力線の近くで最大約 20 μT。IARC 2002 の 2B）。本文は言いすぎないように書いた。
%%   ・送電の電圧の段（2025年の時点の資料で照らし合わせた。二次資料）：日本の送電線の公称電圧は 500・275・220・187・154・110・77・66 kV など、
%%       配電用変電所で 6.6 kV（コトバンク「送電」（日本大百科全書）の検索の抜粋、CGWORLD「電柱のはなし」2019、パナソニック「柱上変圧器」用語解説）。
%%       日本の 500 kV 送電のはじまり：東京電力の房総線、1973年に 500 kV で運用（電気学会「でんきの礎」ishi-2324 の抜粋、IEEE 論文
%%       "500-kV Line Design"（the first 500-kV transmission line in Japan, the Boso Line）の抜粋）。tepco.co.jp・iee.jp・kotobank.jp は作業環境から届かなかった。
%%   ・柱上変圧器 6600 V／210-105 V（単相3線）：JIS C 4304:2013 準拠の製品の仕様（ipros の掲載）、防衛省の入札資料の銘板（愛知電機 6600/105 V 単相3線）。
%%   ・単相3線式（電圧線どうし 200 V、電圧線と中性線 100 V、今のほとんどの住宅で使う）：パナソニック「配線方式」の解説、建築設備用語集（渡辺パイプ）。
%%   ・高圧の契約の力率割引（85 % を上回る 1 % ごとに基本料金を 1 % 割り引く）：関西電力の料金表（biz.kepco.jp の PDF の抜粋）、selectra・enechange の解説。
%%   ・1891年のラウフェン〜フランクフルト間 約175 km の三相交流の送電（IEEE マイルストーン）：ETHW "Milestones: Long Distance Electric Power Transmission
%%       Using Three-Phase Alternating Current, 1891"、IEEE R8 ポーランド支部の記事（2025）。
%%   ・例題の送電線の抵抗 10 Ω、工場・モーターの数は仮の値（本文で「〜とします」と書いた）。
\makeatletter
% 鉄塔の腕の先のがいしと電線  \sxlwire{(x,y)}
\newcommand\sxlwire[1]{\draw[bkInk!70, line width=0.6pt] #1 -- ++(0,-2.4); \filldraw[fill=bkOrange!70, draw=bkInk] ($#1+(0,-2.8)$) circle (0.75);}
% 抵抗（横）  \sxlres{(左)}{(右)}
\newcommand\sxlres[2]{\draw[bkInk, line width=0.8pt, decorate, decoration={zigzag, segment length=1.6mm, amplitude=1.1mm, pre length=1.5mm, post length=1.5mm}] #1 -- #2;}
% 交流の電源の丸  \sxlsrc{(中心)}
\newcommand\sxlsrc[1]{\filldraw[fill=white, draw=bkInk, line width=0.8pt] #1 circle (3); \draw[bkInk, line width=0.7pt] ($#1+(-1.8,0)$) sin ($#1+(-0.9,0.9)$) cos ($#1+(0,0)$) sin ($#1+(0.9,-0.9)$) cos ($#1+(1.8,0)$);}
% フェーザの矢印  \sxlarrow{色}{(始め)}{(終わり)}
\newcommand\sxlarrow[3]{\draw[#1, line width=1.3pt, -{Stealth[length=2.4mm]}] #2 -- #3;}
\tikzset{sxl lab/.style={font=\sffamily\footnotesize, text=bkInk}}
\makeatother

\phchapter{11}{電力を運ぶ3本の線：\\交流電力と三相交流}{送電線の鉄塔では、なぜ電線が3本ずつ組になっている？}{\bklevelrange{2}{5}}{電力を運ぶ3本の線：交流電力と三相交流}
\bkmichishirube{4-9,5-15,6-8,6-9,6-10,電-4,電-8,電-10}{電-12,電-13,電-16}

\begin{goal}
\begin{itemize}[label={\color{bkTheme}●}, leftmargin=1.6\zw, itemsep=0pt, topsep=1pt]
\item \bklevel{2}\enspace 交流の電力が時間とともにゆれること、家の電気と送電線の電気の違いを、図で説明できる。
\item \bklevel{3}\enspace 力率を使って、有効電力・無効電力・皮相電力を計算できる。
\item \bklevel{3}\enspace コンデンサーで力率をよくすると、電線の電流と損失が減るわけを説明できる。
\item \bklevel{4}\enspace 三相交流の Y 結線・$\Delta$ 結線の電圧と電流、三相の電力を計算できる。
\item \bklevel{5}\enspace 三相の電力が時間によらず一定になること、回転磁界ができることを、式で示せる。
\end{itemize}
\end{goal}

\phkoushiki{%
\phk{同じ電力なら、高い電圧で送るほど電線の損失が小さい（損失 $RI^2=R(P/V)^2$）}{高校物理基礎}{ph:電-11:soden}
}

\begin{hatena}
\noindent\begin{minipage}[c]{0.6\linewidth}\raggedright
田んぼや山の中に立つ大きな鉄塔を見上げると、電線は左右に3本ずつ、合わせて6本かかっていて、いちばん上にもう1本細い線があります。家のコンセントの電気は、2本の線で行って帰ってきます。それなら、送電線も2本ずつの組でよさそうです。なぜ3本ずつなのでしょう。
\end{minipage}\hfill
\begin{minipage}[c]{0.37\linewidth}\centering
\begin{tikzpicture}[x=1mm, y=1mm]
  \draw[bkInk, line width=0.9pt] (14,0) -- (18,44) -- (22,0);
  \draw[bkInk!70, line width=0.5pt] (15,10) -- (21,20) -- (15.9,30) -- (20.4,38);
  \draw[bkInk!70, line width=0.5pt] (21,10) -- (15,20) -- (20.1,30) -- (15.6,38);
  \foreach \yy in {38,30,22} {\draw[bkInk, line width=0.9pt] (4,\yy) -- (32,\yy);}
  \foreach \yy in {38,30,22} {\sxlwire{(5,\yy)} \sxlwire{(31,\yy)}}
  \draw[bkInk, line width=0.6pt] (18,44) -- (18,47);
  \filldraw[fill=bkGray!50, draw=bkInk] (18,47) circle (0.6);
  \node[sxl lab, anchor=west] at (19.5,48) {架空地線};
  \node[sxl lab, anchor=north] at (5,17) {1回線};
  \node[sxl lab, anchor=north] at (31,17) {1回線};
  \node[sxl lab] at (18,-3.5) {鉄塔};
\end{tikzpicture}
\end{minipage}
\end{hatena}

\noindent 電-10 では、コイルやコンデンサーをふくむ交流の回路で、電流と電圧の大きさと、山の位置のずれを求めました。この章では、そのずれが\textbf{電力}（1秒あたりに運ぶエネルギー）にどう効くかを調べます。ずれがあると、電流は流れているのに、エネルギーは行ったり来たりするだけ、という部分ができます。そして、発電所から工場や家まで電気を届けるしくみの中心にある\textbf{三相交流}を、数学の 4-9・6-9・6-10 で学んだ「120° ずつずれた3つの波」の物理として読み直します。

\begin{yattemiyou}[3つの波をたしてみよう・家の電気の契約を見よう]
\begin{enumerate}[label={\bkcircled[bkGreen]{\arabic*}}, leftmargin=2.4\zw, itemsep=1pt, topsep=0pt]
\item 表計算のA列に角 $\theta$（0°から30°ごとに360°まで）を入れ、120°ずつずれた3つの波と、その和を作る。\par
{\small\texttt{B2 = SIN(RADIANS(A2))}\quad\texttt{C2 = SIN(RADIANS(A2-120))}\par
\texttt{D2 = SIN(RADIANS(A2-240))}\quad\texttt{E2 = B2+C2+D2}}
\item E列（和）はどうなるか。B列〜D列のグラフもかいてみる。
\item 電気の検針票（電気料金のお知らせ）を家の人といっしょに見て、「契約の電流（たとえば\qty{40}{A}）」と「使った電力量（kWh）」の欄を探す（分電盤にはさわらず、見るだけ）。
\end{enumerate}
\end{yattemiyou}

\noindent やってみようの答え：E列はどの角でも0です（たとえば $\theta=90°$ なら、1と $-0.5$ と $-0.5$ で、和は0）。3つの波は、どの瞬間も、たすと打ち消しあいます。この「和がいつも0」が、鉄塔の電線が3本ずつである理由の中心です（⑦⑧）。検針票の電力量は、この章の\textbf{有効電力}を時間でたしたものです（④）。

\begin{omoidasou}
\bkrecall{4-9}{電力 $p=V_0\sin\omega t\times I_0\sin\omega t$ の平均は $V_0I_0/2$。実効値 $V_{\mathrm e}=V_0/\sqrt2$ を使うと $V_{\mathrm e}I_{\mathrm e}$。電流が60°おくれると平均は半分になり、$\cos$（おくれの角）を力率という。}
\bkrecall{5-15}{$\sin^2$ の1周期の平均は $\tfrac12$（積分で確かめた）。家の\qty{100}{V}は実効値で、山は約\qty{141}{V}。}
\bkrecall{6-10}{交流 $\sqrt2V\sin(\omega t+\theta)$ を、長さ $V$・角 $\theta$ の矢印（フェーザ）で表す。三相交流のフェーザは、120°ずつずれた3本の矢印で、和は0。}
\bkrecall{6-9}{1の3乗根 $1$、$-\tfrac12\pm\tfrac{\sqrt3}{2}j$ の和は0。$e^{j2\pi/3}$ をかけると、矢印が120°回る。}
\bkrecall{電-4}{直流の電力 $P=VI$。抵抗で熱になる電力（ジュール熱）$RI^2$。}
\bkrecall{電-10}{リアクタンス $X_L=\omega L$、$X_C=1/(\omega C)$。抵抗とコイルの直列では $Z=\sqrt{R^2+X^2}$、電流は電圧より $\phi$ だけおくれ、$\tan\phi=X/R$。}
\end{omoidasou}

\noindent この章の記号の約束：瞬時の電圧・電流（その瞬間の値）を小文字の $v$・$i$、実効値を添字のない大文字の $V$・$I$ で書きます（数学の 4-9 の $V_{\mathrm e}$・$I_{\mathrm e}$ と同じ量）。ことわりがなければ、交流の\qty{100}{V}・\qty{10}{A}は実効値です。電流が電圧より何度おくれているかを $\phi$ と書き、電流が進んでいるときは $\phi$ を負にします。

%% ======================== 物理の道具 ========================
\bksection{doguu}{物理の道具}

\phstage{me}{家の電気と送電線を見る}{2}

\bksubsection{1}{家の電気は2本（3本）、送電線は3本}
\label{ph:電-11:sec1}
発電所でつくった電気は、何回も電圧を変えながら、家や工場まで届きます（図\ref{ph:電-11:fig-route}）。電圧を変える機械が\textbf{変圧器}です（しくみは電-12 ⑤）。

\begin{bkfigure}
\begin{tikzpicture}[x=1mm, y=1mm]
  \tikzset{sxl box/.style={bk node, minimum width=22mm, minimum height=10mm, font=\sffamily\footnotesize, align=center}}
  \node[sxl box] (g) at (12,0) {発電所};
  \node[sxl box] (t1) at (60,0) {変電所\\（上げる）};
  \node[sxl box] (t2) at (108,0) {変電所\\（下げる）};
  \node[sxl box] (t3) at (108,-26) {配電用\\変電所};
  \node[sxl box] (p) at (60,-26) {柱上変圧器};
  \node[sxl box] (h) at (12,-26) {家};
  \node[sxl box] (f) at (108,-50) {工場・ビル};
  \draw[bk link] (g) -- (t1);
  \draw[bk link] (t1) -- node[bk link label, above] {三相} node[bk link label, below] {275・500 kV} (t2);
  \draw[bk link] (t2) -- node[bk link label, left, align=right] {三相\\66・154 kV} (t3);
  \draw[bk link] (t3) -- node[bk link label, above] {三相 6.6 kV} (p);
  \draw[bk link] (p) -- node[bk link label, above] {単相3線式} node[bk link label, below] {100・200 V} (h);
  \draw[bk link] (t3) -- node[bk link label, left] {三相 6.6 kV} (f);
\end{tikzpicture}
\bkfigcaption{発電所から家まで（電圧の段は一例。地域や会社で違う）。高い電圧の送電線と、6.6 kV の配電線は三相交流で、柱上変圧器で単相の 100 V・200 V に下げて家に引きこむ。}\label{ph:電-11:fig-route}
\end{bkfigure}

\begin{itemize}[leftmargin=1.2\zw, itemsep=1pt, topsep=2pt, label={\color{\bkTT}●}]
\item \textbf{送電}：日本の送電線は、\qty{500}{kV}・\qty{275}{kV}・\qty{154}{kV}・\qty{66}{kV} などの電圧で、三相交流を送ります（日本で\qty{500}{kV}の送電が始まったのは1973年）。
\item \textbf{配電}：配電用変電所で\qty{6.6}{kV}に下げ、町の電柱の上の3本の線（三相）で配ります。工場や大きなビルは、この\qty{6.6}{kV}を受けとり、自分の変圧器で\qty{200}{V}などの三相に下げて、モーターを回します。
\item \textbf{家}：電柱の上の\textbf{柱上変圧器}が、三相のうちの2本から\qty{6600}{V}を受けとり、\qty{210}{V}と\qty{105}{V}に下げます。家には3本の線で引きこみ（\textbf{単相3線式}）、両はしの2本の間で\qty{200}{V}、まん中の線（中性線）とはしの線の間で\qty{100}{V}を使います（図\ref{ph:電-11:fig-3w}、例題7）。
\end{itemize}
家に届く電圧は、法律（電気事業法施行規則）で、\qty{100}{V}は「\qty{101}{V}の上下\qty{6}{V}をこえない値」（\qtyrange{95}{107}{V}）、\qty{200}{V}は「\qty{202}{V}の上下\qty{20}{V}をこえない値」（\qtyrange{182}{222}{V}）に保つと決められています。柱上変圧器の\qty{105}{V}は、電線での電圧の下がりを見こんで、少し高めにしてあるのです。

\begin{bkfigure}
\begin{tikzpicture}[x=1mm, y=1mm]
  \draw[bkInk, line width=0.9pt, decorate, decoration={coil, aspect=0, segment length=3mm, amplitude=2.4mm}] (8,1) -- (8,29);
  \draw[bkInk, line width=0.8pt] (8,29) -- (8,30) -- (112,30);
  \draw[bkInk, line width=0.8pt] (8,1) -- (8,0) -- (112,0);
  \draw[bkInk, line width=0.8pt] (8,15) -- (68,15);
  \filldraw[bkInk] (8,15) circle (0.7);
  \draw[bkInk, line width=0.7pt] (2,15) -- (8,15) (2,15) -- (2,10);
  \draw[bkInk, line width=0.7pt] (-1,10) -- (5,10) (0,8.8) -- (4,8.8) (1,7.6) -- (3,7.6);
  \node[sxl lab, anchor=east] at (-1.5,13.5) {接地};
  \node[sxl lab, anchor=south west] at (20,30.4) {電圧線};
  \node[sxl lab, anchor=south west] at (20,15.4) {中性線};
  \node[sxl lab, anchor=north west] at (20,-0.4) {電圧線};
  \node[sxl lab, anchor=south, align=center] at (8,31.5) {柱上変圧器};
  \foreach \yb in {15,0} {
    \draw[bkInk, line width=0.8pt] (46,\yb+15) -- (46,\yb+11.5) (46,\yb+3.5) -- (46,\yb);
    \filldraw[fill=bkBlue!12, draw=bkInk, line width=0.8pt] (39,\yb+3.5) rectangle (53,\yb+11.5);
    \node[sxl lab] at (46,\yb+7.5) {100 V};
  }
  \draw[bkInk, line width=0.8pt] (90,30) -- (90,21) (90,9) -- (90,0);
  \filldraw[fill=bkOrange!15, draw=bkInk, line width=0.8pt] (83,9) rectangle (97,21);
  \node[sxl lab, align=center] at (90,15) {200 V};
  \node[sxl lab, anchor=west, align=left] at (55,23) {照明・\\テレビ};
  \node[sxl lab, anchor=west, align=left] at (55,8) {冷蔵庫・\\洗濯機};
  \node[sxl lab, anchor=west, align=left] at (99,15) {IH・\\エアコン};
\end{tikzpicture}
\bkfigcaption{単相3線式。柱上変圧器の二次側のまん中から中性線を出し、地面につなぐ（接地）。中性線と電圧線の間で 100 V、2本の電圧線の間で 200 V を使う。}\label{ph:電-11:fig-3w}
\end{bkfigure}

\textbf{単相}（たんそう）とは、1つの交流（1つの波）を2本の線で送るやり方、\textbf{三相}とは、山の位置が120°ずつずれた3つの交流を、3本の線で送るやり方です。\bkidx{さんそうこうりゅう}{三相交流}単相3線式の3本の線は、1つの波の「プラス側」と「マイナス側」とまん中なので、三相ではありません。

\begin{phmijika}[三相のモーターは身近にある]
\begin{itemize}[leftmargin=1.2\zw, itemsep=1pt, topsep=1pt, label={\color{phMijikaDeep}●}]
\item ビルのエレベーター、工場のポンプやベルトコンベヤー、大きな店の冷凍庫は、三相\qty{200}{V}のモーターで動いています。
\item 学校の大きなエアコンの室外機に、3本の線のつながる「動力」と書いたスイッチがついていることがあります（三相のこと）。
\item 鉄塔の左右の3本ずつは、それぞれ1組の三相（1回線）です。片方が故障しても、もう片方で送り続けられるように2組かけることが多く、いちばん上の線（架空地線）は、雷を受けとめて地面へ流すための線です。
\end{itemize}
\end{phmijika}

\bksubsection{2}{交流の電力は、時間とともにゆれる}
\label{ph:電-11:sec2}
直流なら、電力は電圧×電流で一定です。交流では、電圧も電流も時間とともに変わるので、その瞬間の電力 $p=vi$（\textbf{瞬時電力}）も変わります。図\ref{ph:電-11:fig-pt}に、3つの場合をかきました（電圧の山は $\sqrt2V$、電流の山は $\sqrt2I$。たての目もりは $VI$ を1にしてある）。

\begin{bkfigure}
\begin{tikzpicture}
\begin{axis}[bkplot, axis x line=bottom, axis y line=left, width=60mm, height=44mm, xmin=0, xmax=370, ymin=-1.3, ymax=2.3, xtick={0,90,180,270,360}, ytick={-1,0,1,2}, xlabel={$\omega t$（°）}, title={① 抵抗だけ（$\phi=0°$）}, title style={font=\sffamily\footnotesize}, tick label style={font=\sffamily\footnotesize}]
  \draw[bkGray, line width=0.6pt] (axis cs:0,0) -- (axis cs:370,0);
  \addplot[bkBlue, line width=1.3pt, domain=0:360, samples=121] {2*sin(x)^2};
  \addplot[bkOrange!80!black, line width=0.9pt, dash pattern=on 3pt off 2pt, domain=0:360] {1};
\end{axis}
\end{tikzpicture}\hfill
\begin{tikzpicture}
\begin{axis}[bkplot, axis x line=bottom, axis y line=left, width=60mm, height=44mm, xmin=0, xmax=370, ymin=-1.3, ymax=2.3, xtick={0,90,180,270,360}, ytick={-1,0,1,2}, xlabel={$\omega t$（°）}, title={② コイルだけ（$\phi=90°$）}, title style={font=\sffamily\footnotesize}, tick label style={font=\sffamily\footnotesize}]
  \draw[bkGray, line width=0.6pt] (axis cs:0,0) -- (axis cs:370,0);
  \addplot[bkRed!80!black, line width=1.3pt, domain=0:360, samples=121] {2*sin(x)*sin(x-90)};
\end{axis}
\end{tikzpicture}

\vspace{2mm}
\begin{tikzpicture}
\begin{axis}[bkplot, axis x line=bottom, axis y line=left, width=118mm, height=48mm, xmin=0, xmax=370, ymin=-0.6, ymax=2.1, xtick={0,90,180,270,360}, ytick={0,0.8,1.8}, yticklabels={0,0.8,1.8}, xlabel={$\omega t$（°）}, title={③ 抵抗とコイル（$\cos\phi=0.8$、$\phi\approx37°$）}, title style={font=\sffamily\footnotesize}, tick label style={font=\sffamily\footnotesize}]
  \draw[bkGray, line width=0.6pt] (axis cs:0,0) -- (axis cs:370,0);
  \addplot[bkPurple, line width=1.3pt, domain=0:360, samples=181] {2*sin(x)*sin(x-36.87)};
  \addplot[bkOrange!80!black, line width=0.9pt, dash pattern=on 3pt off 2pt, domain=0:360] {0.8};
  \node[sxl lab, anchor=west, fill=white, inner sep=1pt] at (axis cs:292,1.25) {平均 0.8（破線）};
\end{axis}
\end{tikzpicture}
\bkfigcaption{瞬時電力 $p=vi$ の時間変化（たての目もりは $p/(VI)$）。①はいつもプラスで、平均は $VI$（破線）。②はプラスとマイナスが同じだけで、平均は0。③は少しの間だけマイナスになり、平均は $0.8VI$。どれも、電圧の2倍の速さでゆれる。}\label{ph:電-11:fig-pt}
\end{bkfigure}

\begin{itemize}[leftmargin=1.2\zw, itemsep=1pt, topsep=2pt, label={\color{\bkTT}●}]
\item ①\ 抵抗（電気ストーブ・電気ポット）：電圧と電流の山がそろっているので、$v$ と $i$ はいつも同じ符号で、$p$ はいつもプラスです。電気のエネルギーは、いつも電源から抵抗へ流れ、熱になります。
\item ②\ コイルだけ：電流が90°おくれるので、$p$ がプラスの間（電源からコイルへ）と、マイナスの間（コイルから電源へ）が、ちょうど同じだけあります。コイルは、電流が増える間に磁場にエネルギーをため（$\tfrac12Li^2$：電-8 ⑧）、減る間にそれを電源へ返します。平均すると、エネルギーは何も使われません。
\item ③\ モーターのように抵抗とコイルの両方の性質がある機械：$p$ はほとんどプラスですが、ときどきマイナスになります。平均すると、①より小さくなります。
\end{itemize}
1秒に50回（\qty{50}{Hz}）向きが変わる交流なら、電力は1秒に100回ゆれます。電球や電気ストーブがちらつかないのは、熱く光るフィラメントやヒーターの温度が、そんなに速くは変わらないからです。

\Needspace{16\baselineskip}\phstage{te}{力率で電力を数える}{3}

\bksubsection{3}{瞬時電力と平均の電力：力率}
\label{ph:電-11:sec3}
②の図\ref{ph:電-11:fig-pt}を式にします。電圧と、$\phi$ だけおくれた電流を
\[v=\sqrt2\,V\sin\omega t,\qquad i=\sqrt2\,I\sin(\omega t-\phi)\]
と書きます。瞬時電力は、4-9⑦の積を和に直す式 $2\sin a\sin b=\cos(a-b)-\cos(a+b)$ で、
\pheq{p=vi=2VI\sin\omega t\sin(\omega t-\phi)=VI\cos\phi-VI\cos(2\omega t-\phi)}{W ＝ V × A}
です。右辺の1つめは時間によらない一定の値、2つめは2倍の角周波数 $2\omega$ でゆれる波で、1周期で平均すると0になります（5-15⑨）。だから、平均の電力は
\pheq{P=\overline{p}=VI\cos\phi}{W ＝ V × A}
です。\label{ph:電-11:rikiritsu}この $\cos\phi$ を\textbf{力率}（りきりつ）といいます\bkidx{りきりつ}{力率}（4-9⑧で、電流が60°おくれると半分になった、あの数）。$P$ は、電気がほかの形（熱・光・機械の仕事）に変わって使われる電力なので、\textbf{有効電力}ともいいます。

\begin{phyomi}[$P=VI\cos\phi$]
\begin{itemize}[leftmargin=1.2\zw, itemsep=1pt, topsep=1pt, label={\color{bkTheme}●}]
\item 電圧計と電流計の読み（実効値）をかけた $VI$ に、「電流のうち、電圧と山がそろっている分の割合」$\cos\phi$ をかけたものが、本当に使われる電力。
\item $\phi=0$（抵抗だけ）なら $\cos\phi=1$ で $P=VI$。中学で学んだ「電力＝電圧×電流」（電-4）は、この場合。
\item $\phi=\pm90°$（コイルだけ、コンデンサーだけ）なら $\cos\phi=0$ で $P=0$。電流が流れても、エネルギーは行ったり来たりするだけ。
\end{itemize}
\end{phyomi}

\begin{phbaisuru}[大きく]{電流のおくれ $\phi$}
同じ電圧・同じ電流（$VI$ が同じ）で、$\phi$ を0°・30°・60°・90°と大きくすると、$\cos\phi$ は1・0.87・0.5・0になり、使われる電力は図\ref{ph:電-11:fig-cos}のように減っていきます。おくれが小さいうちは減り方がゆるやかで、90°に近づくと急に減ります。

\begin{bkcenter}
\begin{tikzpicture}
\begin{axis}[bkplot, axis x line=bottom, axis y line=left, width=96mm, height=44mm, xmin=0, xmax=95, ymin=0, ymax=1.12, xtick={0,30,45,60,90}, ytick={0,0.5,0.87,1}, xlabel={電流のおくれ $\phi$（°）}, ylabel={$P/(VI)$}, tick label style={font=\sffamily\footnotesize}]
  \addplot[bkBlue, line width=1.3pt, domain=0:90, samples=91] {cos(x)};
  \fill[bkRed] (axis cs:0,1) circle[radius=1.6pt] (axis cs:30,0.866) circle[radius=1.6pt] (axis cs:60,0.5) circle[radius=1.6pt] (axis cs:90,0) circle[radius=1.6pt];
\end{axis}
\end{tikzpicture}
\end{bkcenter}
\bkfigcaption{同じ $VI$ で、使われる電力（$\cos\phi$）。}\label{ph:電-11:fig-cos}
\end{phbaisuru}

\phboxhead{使われる電力は、抵抗の分だけ}
抵抗 $R$ とコイル（リアクタンス $X$）の直列に電圧 $V$ をかけると、$I=V/Z$、$Z=\sqrt{R^2+X^2}$、力率は $\cos\phi=R/Z$ です（電-10 ⑤のフェーザの直角三角形）。すると
\[P=VI\cos\phi=(ZI)\times I\times\frac{R}{Z}=RI^2\]
となり、平均の電力は、抵抗で熱になる分 $RI^2$ とぴったり同じです。コイルは、エネルギーをためて返すだけで、平均では何も使いません（\bkreidairef{2}）。

\begin{phflow}[交流の電力]
\phflowrow{$p=vi$}{1周期で平均する\\正弦波の電圧・電流}{$P=VI\cos\phi$}
\phflowrow{$P=VI\cos\phi$}{電圧と電流の\\山がそろう（$\phi=0$）}{$P=VI$（中学理科）}
\end{phflow}

\begin{phmijika}[電気製品の力率]
電気ストーブ・電気ポット・アイロンのようなヒーターは、ほとんど抵抗だけなので、力率はほぼ1です。モーターを使う洗濯機・扇風機・冷蔵庫は、コイルの性質のために電流がおくれ、力率が1より小さくなります。電気製品の銘板に「消費電力」（W）と「定格電流」（A）が書いてあれば、$P/(VI)$ で力率の見当がつきます。
\end{phmijika}

\bksubsection{4}{有効電力・無効電力・皮相電力}
\label{ph:電-11:sec4}
電流 $I$ を、電圧と山がそろった分 $I\cos\phi$ と、電圧より90°おくれた分 $I\sin\phi$ の2つに分けて考えます（フェーザの矢印を、電圧の向きと、それに直角な向きに分ける。図\ref{ph:電-11:fig-tri}の左）。

\begin{bkfigure}
\begin{tikzpicture}[x=1mm, y=1mm]
  \sxlarrow{bkBlue}{(0,0)}{(44,0)}
  \node[sxl lab, anchor=south] at (40,1) {電圧 $V$};
  \sxlarrow{bkRed!80!black}{(0,0)}{(32,-24)}
  \node[sxl lab, anchor=north east] at (31,-24.5) {電流 $I$};
  \draw[bkGray, line width=0.6pt, dash pattern=on 2pt off 1.5pt] (32,0) -- (32,-24);
  \draw[bkOrange!80!black, line width=1pt] (0,-0.6) -- (32,-0.6);
  \node[sxl lab, anchor=north west] at (17,-1.4) {$I\cos\phi$};
  \node[sxl lab, anchor=west] at (33,-12) {$I\sin\phi$};
  \draw[bkInk, line width=0.5pt] (7,0) arc (0:-36.87:7);
  \node[sxl lab] at (9.8,-3.3) {$\phi$};
  \begin{scope}[xshift=64mm]
    \draw[bkInk, line width=1.1pt] (0,0) -- (40,0) -- (40,-30) -- cycle;
    \node[sxl lab, anchor=south] at (20,0.8) {有効電力 $P=VI\cos\phi$〔W〕};
    \node[sxl lab, anchor=west, align=left] at (41,-15) {無効電力\\$Q=VI\sin\phi$\\〔var〕};
    \node[sxl lab, anchor=north east, align=right] at (19,-16) {皮相電力\\$S=VI$〔VA〕};
    \draw[bkInk, line width=0.5pt] (9,0) arc (0:-36.87:9);
    \node[sxl lab] at (12.5,-3.6) {$\phi$};
    \draw[bkInk, line width=0.5pt] (37,0) -- (37,-3) -- (40,-3);
  \end{scope}
\end{tikzpicture}
\bkfigcaption{左：電流のフェーザを、電圧にそろった分と90°ずれた分に分ける。右：それぞれに $V$ をかけると、電力の三角形ができる（$S^2=P^2+Q^2$）。}\label{ph:電-11:fig-tri}
\end{bkfigure}

\begin{itemize}[leftmargin=1.2\zw, itemsep=1pt, topsep=2pt, label={\color{\bkTT}●}]
\item \textbf{有効電力}\bkidx{ゆうこうでんりょく}{有効電力} $P=VI\cos\phi$：電圧にそろった電流が運ぶ電力。平均して使われる。単位は W（ワット）。
\item \textbf{無効電力} $Q=VI\sin\phi$：90°ずれた電流が運ぶ電力。電源とコイル・コンデンサーの間を行ったり来たりするだけで、平均すると0。行き来の大きさを表す量で、単位は var（バール。volt-ampere reactive の頭文字）。\bkidx{むこうでんりょく}{無効電力}
\item \textbf{皮相電力}\bkidx{ひそうでんりょく}{皮相電力} $S=VI$：電圧計と電流計の読みをかけた「見かけの電力」。単位は VA（ボルトアンペア）。
\end{itemize}
三角形の3つの辺なので、三平方の定理（$\cos^2\phi+\sin^2\phi=1$）から
\pheq{S^2=P^2+Q^2,\qquad \cos\phi=\frac{P}{S}}{VA$^2$ ＝ W$^2$ ＋ var$^2$}
です。$Q$ は、コイルのように電流がおくれるとき（$\phi>0$）を正、コンデンサーのように電流が進むとき（$\phi<0$）を負と数えます。

\phboxhead{無効電力は、何が「無効」なのか}
③の式を、加法定理（4-9）で $\cos(2\omega t-\phi)=\cos2\omega t\cos\phi+\sin2\omega t\sin\phi$ とほどくと、
\[p=P\,(1-\cos2\omega t)-Q\sin2\omega t\]
と書けます。1つめは、0と $2P$ の間でゆれて、いつも0以上の部分（使われるエネルギーの流れ）。2つめは、山の高さが $Q$ で、プラスとマイナスに同じだけゆれる部分（行き来するだけのエネルギーの流れ）です。$Q$ のエネルギーは使われないので「無効」ですが、その行き来のための電流は電線を本当に流れ、電線を温めます（⑤⑥）。

\phboxhead{変圧器や発電機は VA で大きさを表す}
変圧器や発電機の中の電線が熱くなるのは、流れる電流 $I$ で決まり、力率にはよりません。だから、変圧器の大きさ（容量）は、W ではなく VA（kVA）で書きます。\qty{50}{kVA}の変圧器は、力率1の負荷なら\qty{50}{kW}まで使えますが、力率0.8なら\qty{40}{kW}までです（\bkreidairef{6}）。

\begin{motto}[4]
\textbf{複素数で電力を書く}\par
電圧と電流をフェーザ（6-10⑨）で $V$（角0）、$I$（角 $-\phi$）と書きます。電流の共役な複素数 $\bar I$（6-8⑦。角 $+\phi$）をかけると、大きさは $VI$、角は $\phi$ なので、
\[V\bar I=VI\cos\phi+jVI\sin\phi=P+jQ\]
になります。これを\textbf{複素電力}といいます。共役をとるのは、電流がおくれるとき（コイル）に $Q$ が正になるように決めたためです。たくさんの機械がつながっていても、$P$ どうし、$Q$ どうしをそれぞれたせば、全体の $P+jQ$ が求まります（$S$ どうしは、ふつうたせない）。
\end{motto}

\bksubsection{5}{力率をよくする：コンデンサーを並べる}
\label{ph:電-11:sec5}
工場には、モーターがたくさんあります。モーターの電流はおくれるので、力率が1よりかなり小さくなることがあります（例題3では0.7とした）。同じ仕事（有効電力）をするのに、電線に余分な電流（$I\sin\phi$ の分）が流れ、電線や変圧器で熱が出て、むだになります。

そこで、モーターと並列に\textbf{コンデンサー}をつなぎます。コンデンサーの電流 $I_C=\omega CV$（電-10 ④）は、電圧より90°\textbf{進む}ので、モーターの90°おくれた分を打ち消します（図\ref{ph:電-11:fig-kaizen}）。コンデンサーの無効電力は
\pheq{Q_C=VI_C=\omega CV^2}{var ＝ $\dfrac{1}{\text{s}}$×F×V$^2$}
（負の向き）なので、モーターの無効電力 $Q$ と同じ大きさにすれば、全体の力率が1になります。必要な容量は
\[C=\frac{Q}{\omega V^2}\]
です。工場の受電の設備にある\textbf{進相コンデンサー}が、これです。\bkidx{りきりつかいぜん}{力率改善}

\begin{bkfigure}
\begin{tikzpicture}[x=1mm, y=1mm]
  \sxlarrow{bkBlue}{(0,7)}{(46,7)}
  \node[sxl lab, anchor=south] at (42,7.6) {電圧 $V$};
  \sxlarrow{bkRed!80!black}{(0,0)}{(28,-28.6)}
  \node[sxl lab, anchor=north east] at (27,-29) {モーターの電流};
  \sxlarrow{bkGreen!60!black}{(28,-28.6)}{(28,-0.9)}
  \node[sxl lab, anchor=west, align=left] at (29.5,-15) {コンデンサーの電流\\（90°進む）};
  \draw[bkPurple, line width=1.8pt, -{Stealth[length=2.6mm]}] (0,0) -- (28,0);
  \node[sxl lab, anchor=south] at (13,0.6) {合わせた電流};
  \begin{scope}[xshift=80mm, yshift=-26mm]
    \draw[bkInk, line width=0.8pt] (0,0) -- (0,9) (0,15) -- (0,24) -- (34,24);
    \draw[bkInk, line width=0.8pt] (0,0) -- (34,0);
    \sxlsrc{(0,12)}
    \node[sxl lab, anchor=east] at (-3.6,12) {$V$};
    \filldraw[fill=bkRed!15, draw=bkInk, line width=0.8pt] (16,6) rectangle (24,18);
    \node[sxl lab] at (20,12) {M};
    \draw[bkInk, line width=0.8pt] (20,24) -- (20,18) (20,6) -- (20,0);
    \draw[bkInk, line width=0.8pt] (32,24) -- (32,13.2) (32,10.8) -- (32,0);
    \draw[bkInk, line width=1.2pt] (29,13.2) -- (35,13.2) (29,10.8) -- (35,10.8);
    \node[sxl lab, anchor=west] at (36,12) {$C$};
    \draw[bkPurple, line width=1pt, -{Stealth[length=2mm]}] (4,26.5) -- (12,26.5);
    \node[sxl lab, anchor=west] at (12.5,26.5) {電線の電流};
  \end{scope}
\end{tikzpicture}
\bkfigcaption{力率改善。モーター（M）にコンデンサーを並列につなぐと、コンデンサーの電流（90°進む）がモーターの電流のおくれた分を打ち消し、電線を流れる電流（太い矢印）は電圧にそろって小さくなる。モーターそのものの電流は変わらない。}\label{ph:電-11:fig-kaizen}
\end{bkfigure}

\noindent 大事なのは、モーターそのものの電流や仕事は何も変わらないことです。変わるのは、電源から工場までの電線を流れる電流です。モーターとコンデンサーの間で、無効電力（磁場と電場のエネルギー）を直接やりとりさせ、遠くの発電所まで行き来させないようにする、と考えるとわかりやすいでしょう。

\begin{phbaisuru}[半分]{力率}
同じ有効電力 $P$ を同じ電圧で使うとき、電線の電流は $I=P/(V\cos\phi)$ なので、力率を1から0.5に下げると電流は2倍、電線で熱になる損失 $RI^2$ は4倍になります（図\ref{ph:電-11:fig-pfloss}）。

\begin{bkcenter}
\begin{tikzpicture}
\begin{axis}[bkplot, axis x line=bottom, axis y line=left, width=96mm, height=46mm, xmin=0.45, xmax=1.03, ymin=0, ymax=4.4, xtick={0.5,0.6,0.7,0.8,0.9,1}, ytick={1,2,3,4}, x dir=reverse, xlabel={力率 $\cos\phi$}, ylabel={力率1のときの何倍}, tick label style={font=\sffamily\footnotesize}]
  \addplot[bkBlue, line width=1.3pt, domain=0.5:1, samples=60] {1/x};
  \addplot[bkRed!80!black, line width=1.3pt, dash pattern=on 4pt off 2pt, domain=0.5:1, samples=60] {1/x^2};
  \node[sxl lab, anchor=north west, fill=white, inner sep=1pt] at (axis cs:0.64,1.3) {電流（実線）};
  \node[sxl lab, anchor=west, fill=white, inner sep=1pt] at (axis cs:0.6,3.3) {損失（破線）};
\end{axis}
\end{tikzpicture}
\end{bkcenter}
\bkfigcaption{同じ有効電力を送るときの、力率と電線の電流・損失。力率0.7で損失は約2倍。}\label{ph:電-11:fig-pfloss}
\end{phbaisuru}

\bksubsection{6}{送電の損失と電圧}
\label{ph:電-11:sec6}
発電所から、電力 $P$ を、電圧 $V$・力率 $\cos\phi$ で、往復の抵抗が $R$ の電線で送るとします（まずは単相で考えます。三相も同じ形になる：⑧）。電線の電流と、電線で熱になる損失は
\[I=\frac{P}{V\cos\phi},\qquad P_{\text{損}}=RI^2=\frac{RP^2}{V^2\cos^2\phi}\]
です。\label{ph:電-11:soden}力率が1なら、
\pheq{P_{\text{損}}=R\Bigl(\frac{P}{V}\Bigr)^2}{W ＝ Ω × (W/V)$^2$}
です。送る電力 $P$ と電線（$R$）が決まっているとき、損失は電圧の2乗に反比例します。

\begin{phbaisuru}{送る電圧 $V$}
電圧を2倍にすると、同じ電力を送るのに要る電流は半分、損失は $\tfrac14$ になります。10倍なら、電流は $\tfrac1{10}$、損失は $\tfrac1{100}$ です。\qty{100}{kV}で送ると1割が熱になってしまう電線でも、\qty{500}{kV}にすれば0.4\,\%ですみます（\bkreidairef{5}）。
\end{phbaisuru}

\begin{phflow}[送電の損失]
\phflowrow{$P_{\text{損}}=RI^2$、$P=VI\cos\phi$}{同じ $P$ を送る\\$I$ を消す}{$P_{\text{損}}=\dfrac{RP^2}{V^2\cos^2\phi}$}
\phflowrow{$P_{\text{損}}=\dfrac{RP^2}{V^2\cos^2\phi}$}{力率1}{$P_{\text{損}}=R\Bigl(\dfrac{P}{V}\Bigr)^2$（高校物理基礎）}
\end{phflow}

\phboxhead{電圧は高いほどよいのか}
電圧を上げるほど損失は減りますが、電線と鉄塔の間、電線どうしの間の空気が電気を通さないように（空気の絶縁がこわれないように）、がいし（電線を支える焼き物の部品）を長くし、電線どうしを離し、鉄塔を高くしなければなりません。家の中で何十万ボルトを使うこともできません。そこで、遠くへたくさん送る所だけを高い電圧にし、使う所の近くで何段にも分けて下げています（図\ref{ph:電-11:fig-route}）。電圧を自由に上げ下げできる変圧器（電-12 ⑤）は交流でしか使えないことが、電気を交流で送るようになった大きな理由です。

\begin{phmodel}[交流の電力と送電]
\phignore{電線のコイルの性質（インダクタンス）と、電線と地面の間のコンデンサーの性質（電-13）。電線の温度による抵抗の変化。}
\phassume{電圧・電流は正弦波。負荷は抵抗・コイル・コンデンサーの組み合わせで、電流は電圧に比例する。電線は抵抗 $R$ だけ。}
\phlimit{長い送電線（数百 km）や高い周波数（電-13）。整流器やインバーターのように、電流の波の形がくずれる負荷（電-16）。}
\end{phmodel}

\begin{phmada}[送電線の弱い磁場を長く浴びると、体に影響はある？]
電線を流れる\qty{50}{Hz}・\qty{60}{Hz}の電流は、周りに同じ周波数でゆれる磁場をつくります（電-6）。とても強い磁場は、体の中に電流を起こして、神経を刺激するなどの影響を出すことがわかっています。このため法律（電気設備に関する技術基準を定める省令）は、変電所や送電線の近くの人がいる場所の磁場（磁束密度）の平均を、\qty{200}{\micro T}以下にすると決めています。

一方、ずっと弱い磁場を長い年月浴びることの影響は、はっきりしていません。家の中の磁場の平均は\qtyrange{0.05}{0.1}{\micro T}くらい、電力線の近くでも最大\qty{20}{\micro T}くらいです。\qtyrange{0.3}{0.4}{\micro T}以上の磁場を長く浴びる子どもで、白血病が少し多いという疫学調査（たくさんの人の統計）の結果がいくつかあります。世界保健機関（WHO）の環境保健クライテリア（2007年）は、調べ方のかたより（選ばれた人のかたより、浴びた量の見積もりの誤り）が残ることや、動物や細胞の実験で、弱い磁場が体の働きを変えるしくみが確かめられていないことから、「因果関係があると考えるほどには証拠は強くないが、関心を残すには十分に強い」とまとめています。統計に見える弱い関連が、本当に磁場のせいなのか、調べ方のかたよりのせいなのかは、まだわかっていません。
\end{phmada}

\phkokomade

\Needspace{16\baselineskip}\phstage{shiki}{三相交流を式で読む}{5}

\bksubsection{7}{三相交流：Y 結線と $\Delta$ 結線}
\label{ph:電-11:sec7}
発電所の発電機には、120°ずつずらして3つのコイルが巻いてあり、磁石（電磁石）が回ると、3つのコイルに山の位置が120°ずつずれた起電力が生まれます（電-8 ⑤の発電機を3つ重ねたもの。電-12 ⑧）。3つの相を a・b・c とよび、起電力を
\[e_{\mathrm a}=\sqrt2E\sin\omega t,\quad e_{\mathrm b}=\sqrt2E\sin(\omega t-120°),\quad e_{\mathrm c}=\sqrt2E\sin(\omega t-240°)\]
と書きます（$E$ は実効値）。やってみようで確かめたとおり、どの瞬間も $e_{\mathrm a}+e_{\mathrm b}+e_{\mathrm c}=0$ です（4-9 の三相交流、6-9 の1の3乗根、6-10 のフェーザの正三角形）。

\phboxhead{Y 結線：帰りの線が要らない}
3つのコイルの片方のはしを1点（\textbf{中性点}）にまとめ、もう片方のはしから線を出すつなぎ方を、形から \textbf{Y 結線}（スター結線）といいます。\bkidx{わいけっせん}{Y 結線}負荷も同じ形にして、3つとも同じインピーダンス（\textbf{平衡}）にすると、3つの電流も大きさが同じで、120°ずつずれます。すると中性点どうしを結ぶ線（中性線）の電流は
\[i_{\mathrm a}+i_{\mathrm b}+i_{\mathrm c}=0\]
で、いつも0です。電流が流れない線は、なくしてもかまいません。単相3つなら行きと帰りで6本の線が要るところを、三相なら3本ですませられるのです（図\ref{ph:電-11:fig-Y}）。

\begin{bkfigure}
\begin{tikzpicture}[x=1mm, y=1mm]
  \foreach \ang in {90,210,330} {
    \draw[bkInk, line width=0.8pt] (0,0) -- (\ang:10);
    \sxlsrc{(\ang:13)}
    \draw[bkInk, line width=0.8pt] (\ang:16) -- (\ang:19);
    \draw[bkInk, line width=0.8pt] (66,0) -- ($(66,0)+(\ang:3)$);
    \sxlres{($(66,0)+(\ang:3)$)}{($(66,0)+(\ang:16)$)}
    \draw[bkInk, line width=0.8pt] ($(66,0)+(\ang:16)$) -- ($(66,0)+(\ang:19)$);
  }
  \node[sxl lab, anchor=east] at (-3.5,13) {$e_{\mathrm a}$};
  \node[sxl lab, anchor=south east] at (-14,-2.5) {$e_{\mathrm b}$};
  \node[sxl lab, anchor=south west] at (14,-2.5) {$e_{\mathrm c}$};
  \filldraw[bkInk] (0,0) circle (0.8);
  \filldraw[bkInk] (66,0) circle (0.8);
  \node[sxl lab, anchor=south east] at (-1.2,1.2) {中性点};
  \node[sxl lab, anchor=south west] at (67.5,1.8) {中性点};
  \draw[bkInk, line width=0.8pt] (0,19) -- (66,19);
  \draw[bkInk, line width=0.8pt] (330:19) -- ($(66,0)+(210:19)$);
  \draw[bkInk, line width=0.8pt] (210:19) -- ++(0,-7) -- ($(66,0)+(330:19)+(5,-7)$) -- ++(0,7) -- ($(66,0)+(330:19)$);
  \draw[bkInk!60, line width=0.7pt, dash pattern=on 2pt off 1.5pt] (0,0) -- (30:9) -- ($(66,0)+(150:9)$) -- (66,0);
  \node[sxl lab, anchor=south] at (33,5.2) {中性線（電流は0なので、なくせる）};
  \node[sxl lab, anchor=south] at (33,19.5) {線 a};
  \node[sxl lab, anchor=south] at (33,-9) {線 c};
  \node[sxl lab, anchor=north] at (33,-17) {線 b};
  \node[sxl lab, anchor=west] at (76,14) {負荷};
\end{tikzpicture}
\bkfigcaption{Y 結線の電源と負荷。3つの負荷が同じなら、中性線（破線）には電流が流れない。}\label{ph:電-11:fig-Y}
\end{bkfigure}

\phboxhead{線と線の間の電圧は $\sqrt3$ 倍}
Y 結線では、1つのコイルの電圧（\textbf{相電圧} $V_{\mathrm p}$）と、2本の線の間の電圧（\textbf{線間電圧} $V$）が違います。線 a と線 b の間の電圧は $e_{\mathrm a}-e_{\mathrm b}$ です。フェーザで書くと、長さ $E$ の2本の矢印（角0°と $-120°$）の差です（図\ref{ph:電-11:fig-sqrt3}）。

\begin{bkfigure}
\begin{tikzpicture}[x=1mm, y=1mm]
  \sxlarrow{bkBlue}{(0,0)}{(30,0)}
  \node[sxl lab, anchor=south] at (24,0.6) {$e_{\mathrm a}$};
  \sxlarrow{bkGreen!60!black}{(0,0)}{(-15,-25.98)}
  \node[sxl lab, anchor=east] at (-16,-25) {$e_{\mathrm b}$};
  \sxlarrow{bkOrange!80!black}{(0,0)}{(-15,25.98)}
  \node[sxl lab, anchor=east] at (-16,25) {$e_{\mathrm c}$};
  \sxlarrow{bkRed!80!black}{(-15,-25.98)}{(30,0)}
  \node[sxl lab, anchor=north west] at (9,-14) {$e_{\mathrm a}-e_{\mathrm b}$（長さ $\sqrt3E$）};
  \draw[bkInk, line width=0.5pt] (6,0) arc (0:-120:6);
  \node[sxl lab] at (5.5,-6.8) {120°};
  \draw[bkInk, line width=0.5pt] (22,0) arc (180:210:8);
  \node[sxl lab] at (17,-2.6) {30°};
  \begin{scope}[xshift=72mm]
    \draw[bkInk, line width=0.9pt] (0,-20) -- (40,-20) -- (20,14.64) -- cycle;
    \draw[bkGray, line width=0.6pt, dash pattern=on 2pt off 1.5pt] (20,14.64) -- (20,-20);
    \node[sxl lab, anchor=north] at (20,-20.6) {$\sqrt3E$};
    \node[sxl lab, anchor=south east] at (9,-3) {$E$};
    \node[sxl lab, anchor=south west] at (31,-3) {$E$};
    \node[sxl lab, align=center, fill=white, inner sep=0.8pt] at (20,5) {120°};
    \node[sxl lab, anchor=west, align=left] at (21,-12) {半分は\\$E\cos30°$};
  \end{scope}
\end{tikzpicture}
\bkfigcaption{線間電圧のフェーザ。左：$e_{\mathrm a}-e_{\mathrm b}$ は、$e_{\mathrm b}$ の先から $e_{\mathrm a}$ の先への矢印で、$e_{\mathrm a}$ より30°進む。右：頂角120°の二等辺三角形の底辺は $2E\cos30°=\sqrt3E$。}\label{ph:電-11:fig-sqrt3}
\end{bkfigure}

\noindent 頂角が120°、2辺が $E$ の二等辺三角形の底辺なので、長さは $2E\cos30°=\sqrt3E$ です。複素数でも、$E-E\bigl(-\tfrac12-\tfrac{\sqrt3}{2}j\bigr)=E\bigl(\tfrac32+\tfrac{\sqrt3}{2}j\bigr)$、大きさ $E\sqrt{\tfrac94+\tfrac34}=\sqrt3E$ と確かめられます。だから、
\pheq{V=\sqrt3\,V_{\mathrm p},\qquad I=I_{\mathrm p}\qquad\text{（Y 結線）}}{V ＝ V}
（線を流れる\textbf{線電流} $I$ は、コイルを流れる\textbf{相電流} $I_{\mathrm p}$ と同じ）です。工場の「三相\qty{200}{V}」は線間電圧で、1相の電圧は $200\div\sqrt3\approx\qty{115}{V}$ です（\bkreidairef{8}）。

\phboxhead{$\Delta$ 結線：線の電流が $\sqrt3$ 倍}
3つのコイル（負荷）を三角形の輪につなぎ、3つの角から線を出すつなぎ方を \textbf{$\Delta$ 結線}（デルタ結線）といいます\bkidx{でるたけっせん}{$\Delta$ 結線}（図\ref{ph:電-11:fig-delta}）。1つの負荷には2本の線の間の電圧がそのままかかるので $V_{\mathrm p}=V$ です。かわりに、1本の線には、となりあう2つの負荷の電流の差が流れるので、Y の電圧と同じ考えで
\pheq{V=V_{\mathrm p},\qquad I=\sqrt3\,I_{\mathrm p}\qquad\text{（$\Delta$ 結線）}}{A ＝ A}

\begin{bkfigure}
\begin{tikzpicture}[x=1mm, y=1mm]
  \coordinate (A) at (60,30); \coordinate (B) at (45,4); \coordinate (C) at (75,4);
  \sxlres{(A)}{(B)} \sxlres{(B)}{(C)} \sxlres{(C)}{(A)}
  \foreach \pp in {A,B,C} {\filldraw[bkInk] (\pp) circle (0.7);}
  \draw[bkInk, line width=0.8pt] (A) -- (0,30);
  \draw[bkInk, line width=0.8pt] (B) -- (0,4);
  \draw[bkInk, line width=0.8pt] (C) -- (75,-4) -- (0,-4);
  \node[sxl lab, anchor=south] at (10,30.4) {線 a};
  \node[sxl lab, anchor=south] at (10,4.4) {線 b};
  \node[sxl lab, anchor=north] at (10,-4.4) {線 c};
  \draw[bkRed!80!black, line width=1pt, -{Stealth[length=2mm]}] (24,33) -- (36,33);
  \node[sxl lab, anchor=south] at (30,33.4) {線電流 $I$};
  \draw[bkBlue, line width=1pt, -{Stealth[length=2mm]}] (71.5,25) -- (75.5,18);
  \node[sxl lab, anchor=west] at (77,22) {相電流 $I_{\mathrm p}$};
  \draw[bkBlue, line width=1pt, -{Stealth[length=2mm]}] (48.5,25) -- (44.5,18);
  \node[sxl lab, anchor=east] at (43,22) {相電流 $I_{\mathrm p}$};
  \node[sxl lab, anchor=west, align=left] at (84,4) {1つの負荷に\\線間電圧 $V$};
\end{tikzpicture}
\bkfigcaption{$\Delta$ 結線の負荷。線 a の電流は、角 a から出る2つの負荷の電流（相電流）の差で、大きさは相電流の $\sqrt3$ 倍。}\label{ph:電-11:fig-delta}
\end{bkfigure}
になります。平衡した負荷なら、$\Delta$ の1相のインピーダンス $Z_\Delta$ と、Y の1相 $Z_{\mathrm Y}$ は、$Z_{\mathrm Y}=Z_\Delta/3$ のとき、線から見て同じはたらきをします（線電流を2通りに計算して等しくおく：\bkreidairef{9}）。

\begin{phyomi}[$V=\sqrt3V_{\mathrm p}$]
\begin{itemize}[leftmargin=1.2\zw, itemsep=1pt, topsep=1pt, label={\color{bkTheme}●}]
\item 2つの相の電圧は同じ大きさでも、120°ずれているので、差は2倍にはならず、$\sqrt3\approx1.73$ 倍になる。
\item 単相の「プラス側とマイナス側」（180°ずれ）なら差は2倍。単相3線式の 100 V と 200 V がこれ。
\end{itemize}
\end{phyomi}

\bksubsection{8}{三相の電力と、時間によらない電力}
\label{ph:電-11:sec8}
三相の負荷の電力は、3つの相の電力の和です。平衡していれば3つとも同じなので、
\[P=3V_{\mathrm p}I_{\mathrm p}\cos\phi\]
です。Y 結線なら $V_{\mathrm p}=V/\sqrt3$、$I_{\mathrm p}=I$、$\Delta$ 結線なら $V_{\mathrm p}=V$、$I_{\mathrm p}=I/\sqrt3$ なので、どちらでも
\pheq{P=\sqrt3\,VI\cos\phi}{W ＝ V × A}
になります（$V$・$I$ は線間電圧と線電流。$\phi$ は1相の電圧と電流のずれ）。線ではかれる量だけで書けるので、こちらがよく使われます。

\phboxhead{3つの瞬時電力をたすと、一定になる}
③の式を、3つの相に使います。相 b・c では、電圧も電流も $120°$・$240°$ ずつずれているので、2倍の角でゆれる部分は $240°$・$480°$ ずれます。
\[p=3V_{\mathrm p}I_{\mathrm p}\cos\phi-V_{\mathrm p}I_{\mathrm p}\bigl\{\cos(2\omega t-\phi)+\cos(2\omega t-\phi-240°)+\cos(2\omega t-\phi-480°)\bigr\}\]
$480°$ は1周（360°）をひくと $120°$ なので、かっこの中は、120°ずつずれた3つの波の和で、0です（やってみようと同じ）。だから
\pheq{p=3V_{\mathrm p}I_{\mathrm p}\cos\phi=P\qquad\text{（どの時刻でも一定）}}{W ＝ W}
です。単相の電力は1秒に100回（\qty{50}{Hz}なら）ゆれましたが、平衡した三相の電力は、まったくゆれません（図\ref{ph:電-11:fig-3p}、\bkreidairef{11}）。

\begin{bkfigure}
\begin{tikzpicture}
\begin{axis}[bkplot, axis x line=bottom, axis y line=left, width=118mm, height=56mm, xmin=0, xmax=370, ymin=-0.4, ymax=2.9, xtick={0,90,180,270,360}, ytick={0,1,2,2.55}, yticklabels={0,1,2,2.55}, xlabel={$\omega t$（°）}, ylabel={$p/(V_{\mathrm p}I_{\mathrm p})$}, tick label style={font=\sffamily\footnotesize}]
  \draw[bkGray, line width=0.6pt] (axis cs:0,0) -- (axis cs:370,0);
  \addplot[bkBlue, line width=1pt, domain=0:360, samples=181] {2*sin(x)*sin(x-31.79)};
  \addplot[bkGreen!60!black, line width=1pt, dash pattern=on 3pt off 1.5pt, domain=0:360, samples=181] {2*sin(x-120)*sin(x-120-31.79)};
  \addplot[bkOrange!80!black, line width=1pt, dash pattern=on 1.2pt off 1.2pt, domain=0:360, samples=181] {2*sin(x-240)*sin(x-240-31.79)};
  \addplot[bkRed!80!black, line width=1.8pt, domain=0:360] {2.55};
  \node[sxl lab, anchor=north, fill=white, inner sep=1pt] at (axis cs:190,2.47) {3つの和（太い線）＝一定};
\end{axis}
\end{tikzpicture}
\bkfigcaption{三相の瞬時電力（力率0.85、相電圧・相電流の実効値を1にした）。3つの相の電力（実線・破線・点線）はそれぞれ $-0.15$ と1.85の間でゆれるが、和はいつも $3\times0.85=2.55$。}\label{ph:電-11:fig-3p}
\end{bkfigure}

\noindent 電力が一定ということは、三相の発電機を回すのに要るトルクも、三相のモーターが出すトルクも、時間でゆれないということです。単相のモーターは、1秒に100回ゆれる電力のために、ぶるぶるとふるえやすくなります。大きなモーターや発電機が三相でつくられる理由の1つです。

\phboxhead{はてな？の答え}
\bkhatenaref の答え：送電線は三相交流を送っているからです。三相交流の3本の電流は、平衡していれば和がいつも0なので、帰りの線が要りません。単相なら2本で送る電力を、三相は3本で送りますが、同じ電力・同じ電圧・同じ損失で比べると、使う銅は単相の $\tfrac34$ ですみます（\bkreidairef{12}）。しかも、三相の電力は時間でゆれず一定で、3つのコイルに流すとぐるぐる回る磁場（⑨）ができて、モーターを回せます。鉄塔の左右の3本ずつは、それぞれ1組の三相（1回線）で、片方が止まっても送り続けられるように2回線をかけています。

\begin{phflow}[三相の電力]
\phflowrow{$p=\displaystyle\sum_{\text{3相}}v_ki_k$}{平衡した三相\\正弦波}{$p=3V_{\mathrm p}I_{\mathrm p}\cos\phi$（一定）}
\phflowrow{$P=3V_{\mathrm p}I_{\mathrm p}\cos\phi$}{Y でも $\Delta$ でも\\線の量で書く}{$P=\sqrt3VI\cos\phi$}
\end{phflow}

\begin{column}{三相交流で、遠くへ送った}
1891年、ドイツのラウフェンにある水力発電所から、約\qty{175}{km}離れたフランクフルトの電気の博覧会の会場まで、三相交流で電気を送る実験が成功しました。当時は、直流で送るか交流で送るかの議論が続いていましたが、高い電圧で送って変圧器で下げる交流と、三相のモーターの便利さが広く知られ、この後、世界の送電は三相交流が中心になりました。この実験は、電気の技術の歴史に残る出来事として、アメリカの電気電子学会（IEEE）の「マイルストーン」に選ばれています。
\end{column}

\bksubsection{9}{回転磁界：3つのコイルで、回る磁場をつくる}
\label{ph:電-11:sec9}
円の周りに、向きを120°ずつずらして3つのコイルを置き、三相交流を流します（図\ref{ph:電-11:fig-rot}）。コイル a は0°の向き、コイル b は120°の向き、コイル c は240°の向きに、それぞれの電流に比例した磁場をつくります。円の中心の磁場の大きさを、電流が山のときに $B_0$ とすると、3つの磁場は
\[\text{a：}B_0\cos\omega t\ \text{（0°の向き）},\qquad \text{b：}B_0\cos(\omega t-120°)\ \text{（120°の向き）},\]
\[\text{c：}B_0\cos(\omega t-240°)\ \text{（240°の向き）}\]
です（電流を $\cos$ で書いた。$\sin$ と同じ形で、始まりの時刻がずれるだけ）。中心の磁場は、この3つの矢印のベクトルの和です。

\phboxhead{$\omega t=0°$ のとき}
a は $B_0$（0°の向き）、b と c は $\cos(-120°)=\cos(-240°)=-\tfrac12$ なので、それぞれの向きと反対に $\tfrac12B_0$ です。120°の向きと反対は $-60°$、240°の向きと反対は $+60°$ の向きなので、2つのたての成分は打ち消しあい、横の成分は $\tfrac12B_0\cos60°\times2=\tfrac12B_0$。全体は0°の向きに $B_0+\tfrac12B_0=\tfrac32B_0$ です。

\phboxhead{$\omega t=90°$ のとき}
a は $\cos90°=0$。b は $\cos(-30°)=\tfrac{\sqrt3}{2}$ で120°の向き、c は $\cos(-150°)=-\tfrac{\sqrt3}{2}$ で、240°の向きの反対の60°の向きです。横の成分は打ち消しあい、たての成分は $\tfrac{\sqrt3}{2}B_0\sin60°\times2=\tfrac32B_0$。全体は90°の向きに $\tfrac32B_0$ です。

\noindent 大きさは同じ $\tfrac32B_0$ のまま、向きだけが $\omega t$ と同じ角だけ回っています。Python で30°ごとに計算しても、いつも大きさ1.5、向きは $\omega t$ と同じでした。

\begin{bkfigure}
\begin{tikzpicture}[x=1mm, y=1mm]
  \foreach \tt/\xs in {0/0,60/45,120/90} {
    \begin{scope}[xshift=\xs mm]
      \draw[bkGray, line width=0.5pt] (0,0) circle (17);
      \draw[bkRed!80!black, line width=2.2pt, -{Stealth[length=3mm]}] (0,0) -- ({15*cos(\tt)},{15*sin(\tt)});
      \foreach \ang/\cc in {0/bkBlue,120/bkGreen!60!black,240/bkOrange!80!black} {
        \pgfmathsetmacro\ss{10*cos(\tt-\ang)}
        \draw[\cc, line width=1.1pt, -{Stealth[length=2mm]}] (0,0) -- ({\ss*cos(\ang)},{\ss*sin(\ang)});
        \filldraw[fill=\cc!25, draw=\cc] ({19*cos(\ang)},{19*sin(\ang)}) circle (1.6);
      }
      \node[sxl lab] at (0,-23.5) {$\omega t=\tt°$};
    \end{scope}
  }
  \node[sxl lab, anchor=north] at (19,-2.2) {a};
  \node[sxl lab, anchor=east] at (-11.5,17.5) {b};
  \node[sxl lab, anchor=east] at (-11.5,-17.5) {c};
\end{tikzpicture}
\bkfigcaption{回転磁界。外の丸がコイル a・b・c の置き場所。細い矢印が3つのコイルの磁場、太い矢印がその和。和の大きさは変わらず、向きが $\omega t$ と同じだけ回る。}\label{ph:電-11:fig-rot}
\end{bkfigure}

\phboxhead{複素数で、どの時刻でも確かめる}
向き $\theta$ の長さ1の矢印を $e^{j\theta}$（6-10）で表すと、中心の磁場は
\[B_0\sum_{k=0}^{2}\cos\Bigl(\omega t-\frac{2\pi k}{3}\Bigr)e^{j2\pi k/3}\]
です。$\cos x=\tfrac12(e^{jx}+e^{-jx})$（6-10⑧）でほどくと、1つめの部分は $\tfrac12e^{j\omega t}$ が3つで $\tfrac32e^{j\omega t}$、2つめの部分は $\tfrac12e^{-j\omega t}e^{j4\pi k/3}$ の和で、$e^{j4\pi k/3}$（$k=0,1,2$）は1の3乗根の3つなので、和は0です（6-9）。だから
\pheq{\text{中心の磁場}=\frac32B_0\,e^{j\omega t}}{T ＝ T}
で、大きさ $\tfrac32B_0$ の磁場が、角速度 $\omega$ で回ります。これを\textbf{回転磁界}といいます\bkidx{かいてんじかい}{回転磁界}（6-9 の「三相交流とモーター」の約束）。この中に方位磁針や磁石を置くと、磁場に引かれていっしょに回ります。これが、三相のモーター（電-12 ⑦⑧）の原理です。

\phboxhead{回る速さと極の数}
図\ref{ph:電-11:fig-rot}のコイルの組は、N極とS極を1つずつもつ磁場（\textbf{2極}）をつくり、交流の1周期で1回転します。\qty{50}{Hz}なら1秒に50回転、1分に3000回転です。コイルの組を円の周りに2回繰り返して置くと（\textbf{4極}）、1周期で半回転しかしないので、1分に1500回転になります。一般に、磁場の回る速さ（\textbf{同期速度}）は
\pheq{n_{\mathrm s}=\frac{120f}{\text{極の数}}\ \ \text{〔min$^{-1}$〕}}{min$^{-1}$ ＝ Hz × s/min}
です（$60f$ を、極の対の数でわる。\bkreidairef{13}）。

\bksubsection{10}{つり合わない三相と、対称座標法}
\label{ph:電-11:sec10}
家々の単相の負荷は、三相の3つの相に分けてつなぎますが、いつも同じだけ使われるわけではありません。3つの相の電流が等しくない（\textbf{不平衡}）と、中性線に電流が流れ（チャレンジ問題）、モーターには逆向きに回ろうとする磁場がまざって、熱くなります。

不平衡な3つの電流 $I_{\mathrm a}$・$I_{\mathrm b}$・$I_{\mathrm c}$（フェーザ）は、120°回す数 $a=e^{j2\pi/3}$（6-9 の1の3乗根）を使って、つり合った3つの組の和に分けられます。
\[I_0=\frac{I_{\mathrm a}+I_{\mathrm b}+I_{\mathrm c}}{3},\quad I_1=\frac{I_{\mathrm a}+aI_{\mathrm b}+a^2I_{\mathrm c}}{3},\quad I_2=\frac{I_{\mathrm a}+a^2I_{\mathrm b}+aI_{\mathrm c}}{3}\]
$I_1$ は、ふつうの順（a→b→c）に回る三相（\textbf{正相}）、$I_2$ は逆の順に回る三相（\textbf{逆相}）、$I_0$ は3つの相に同じ向きに流れる分（\textbf{零相}）の大きさです。これを\textbf{対称座標法}といいます。\bkidx{たいしょうざひょうほう}{対称座標法}送電線が雷で地面とつながる故障（地絡）のときの電流の計算や、逆相の電流で熱くなるモーターを守るしかけに使われます。詳しい計算は電力系統の専門の本で学びます。

\begin{phsaki}{電力系統工学}
この章では、1つの電源と1つの負荷を考えました。本物の電力系統では、何百もの発電所と、数えきれない負荷が網の目のようにつながっています。どの線にどれだけの電力が流れるか（電力潮流）、故障のときにどれだけの電流が流れるか（故障計算。対称座標法を使う）、発電と使う量のつり合いをどう保つか（周波数の制御：電-16）を計算する学問が、\textbf{電力系統工学}です。高い電圧の直流で送る送電（直流送電）や、半導体のスイッチで電力の流れを自由に変える装置も使われています（電-16）。
\end{phsaki}

\begin{senmonsho}[電力工学の教科書ふうの書き方]
\begin{bkhonbun}
平衡三相回路において，線間電圧を $V$〔V〕，線電流を $I$〔A〕，負荷の力率を $\cos\theta$ とすると，三相電力は
\[P=\sqrt3VI\cos\theta\ \text{〔W〕},\qquad Q=\sqrt3VI\sin\theta\ \text{〔var〕}\]
で表される．また，三相の瞬時電力の和は時間に無関係に一定である．
\end{bkhonbun}
\bkyaku{電気の専門書は、電流のおくれの角を $\theta$ と書くことが多い（この本の $\phi$）。$V$・$I$ は線間電圧と線電流の実効値。$Q$ の式は、1相の $V_{\mathrm p}I_{\mathrm p}\sin\phi$ の3倍を、$P$ と同じように線の量で書き直したもの。}
\end{senmonsho}

\begin{taisetsu}
\def\tsl#1{\makebox[6.6\zw][l]{\sffamily\bfseries\small\color{\bkTT}#1}}
\tsl{瞬時電力}$p=vi=VI\cos\phi-VI\cos(2\omega t-\phi)$\\[3pt]
\tsl{3つの電力}$P=VI\cos\phi$〔W〕、$Q=VI\sin\phi$〔var〕、$S=VI$〔VA〕、$S^2=P^2+Q^2$\\[3pt]
\tsl{力率改善}$Q_C=\omega CV^2$、$C=\dfrac{Q}{\omega V^2}$\\[3pt]
\tsl{送電の損失}$P_{\text{損}}=\dfrac{RP^2}{V^2\cos^2\phi}$（電圧2倍で $\tfrac14$）\\[3pt]
\tsl{三相}Y：$V=\sqrt3V_{\mathrm p}$、$\Delta$：$I=\sqrt3I_{\mathrm p}$、$P=\sqrt3VI\cos\phi$（瞬時電力も一定）\\[3pt]
\tsl{回転磁界}大きさ $\tfrac32B_0$、同期速度 $n_{\mathrm s}=\dfrac{120f}{\text{極の数}}$〔min$^{-1}$〕
\end{taisetsu}

\begin{kotoba}
\small
\bkword{瞬時電力}{その瞬間の電圧×電流 $p=vi$。交流では電圧の2倍の速さでゆれる。}
\bkword{力率}{有効電力÷皮相電力 $\cos\phi$。電流のうち、電圧と山がそろった分の割合。}
\bkword{有効電力・無効電力・皮相電力}{平均して使われる電力〔W〕、行き来するだけの電力の大きさ〔var〕、電圧×電流〔VA〕。}
\bkword{力率改善}{コンデンサーを並列につないで、無効電力を打ち消し、電線の電流を減らすこと。}
\bkword{三相交流}{山の位置が120°ずつずれた3つの交流。3本の線で送れる。}
\bkword{Y 結線・$\Delta$ 結線}{3つの相を1点にまとめるつなぎ方と、三角形の輪にするつなぎ方。}
\bkword{回転磁界}{三相交流を3つのコイルに流してできる、一定の大きさで回る磁場。}
\end{kotoba}

\bkkigocards{\bkkigocard{$Q$}{キュー}{無効電力〔var〕。電-11 では電荷を記号に使わない。電-10 の Q 値、第1巻の体積流量とは別}{reactive power}{電-11}}{\bkkigocard{$S$}{エス}{皮相電力〔VA〕。電-11 では面積を記号に使わない。第2巻のエントロピー $S$ とは別}{apparent power}{電-11}}
\bkkigocards{\bkkigocard{$V_{\mathrm p}$}{ブイ・ピー}{三相の1相の電圧（相電圧）。相電流は $I_{\mathrm p}$。添字なしは線間電圧・線電流}{phase voltage}{電-11}}{\bkkigocard{$n_{\mathrm s}$}{エヌ・エス}{同期速度（回転磁界が1分に回る数）〔min$^{-1}$〕。$120f$÷極の数}{synchronous speed}{電-11}}

%% ======================== 例題 ========================
\Needspace{10\baselineskip}\bksection{reidai}{例題でたしかめよう}

\label{ph:電-11:ex1}
\begin{phreidai}[電気製品の3つの電力]
\qty{100}{V}の交流につないだ機械に\qty{10}{A}の電流が流れ、力率は0.8（電流がおくれる）でした。有効電力・無効電力・皮相電力を求めましょう。
\phbamen
モーターをふくむ機械の、電圧計・電流計の読みと力率から、電力を求める場面です。
\phmushi
電圧・電流は正弦波とします（波の形のくずれを無視）。
\phzu
$V=\qty{100}{V}$、$I=\qty{10}{A}$、$\cos\phi=0.8$。$\sin\phi=\sqrt{1-0.8^2}=0.6$。
\phriyu
④の3つの電力の式は、実効値と力率だけで書けるからです。
\phtateru
\[P=VI\cos\phi,\qquad Q=VI\sin\phi,\qquad S=VI\]
\phtoku
$P=100\times10\times0.8=\qty{800}{W}$、$Q=100\times10\times0.6=\qty{600}{var}$、$S=100\times10=\qty{1000}{VA}$。
\phtashikame
$800^2+600^2=\num{1000000}=1000^2$ で、$S^2=P^2+Q^2$ が成り立つ。$P$ は $S$ より小さい（力率が1より小さいから）。
\kotae 有効電力\qty{800}{W}、無効電力\qty{600}{var}、皮相電力\qty{1000}{VA}
\end{phreidai}
\begin{jibun}
\qty{200}{V}につないだ機械に\qty{5}{A}流れ、力率は0.6でした。有効電力・無効電力・皮相電力を求めましょう。
\bkjans{$P=\qty{600}{W}$、$Q=\qty{800}{var}$、$S=\qty{1000}{VA}$}
\end{jibun}

\label{ph:電-11:ex2}
\begin{phreidai}[抵抗とコイルの直列で使われる電力]
抵抗\qty{4}{\ohm}とコイル（リアクタンス\qty{3}{\ohm}）の直列に、\qty{100}{V}の交流をつなぎました（電-10 例題5 と同じ回路）。電流、力率、使われる電力を求め、抵抗で熱になる電力と等しいことを確かめましょう。
\phbamen
高校の入試でよく見る形の回路で、交流の電力を求める場面です。
\phmushi
コイルの導線の抵抗（\qty{4}{\ohm}にふくめる）、コイルの鉄心の損失を無視します。
\phzu
$R=\qty{4}{\ohm}$、$X=\qty{3}{\ohm}$、$V=\qty{100}{V}$。
\phriyu
電流は $V/Z$（電-10 ⑤）、力率は $R/Z$、電力は③の $VI\cos\phi$ です。
\phtateru
\[Z=\sqrt{R^2+X^2},\qquad I=\frac{V}{Z},\qquad \cos\phi=\frac{R}{Z},\qquad P=VI\cos\phi\]
\phtoku
$Z=\sqrt{16+9}=\qty{5}{\ohm}$、$I=100\div5=\qty{20}{A}$、$\cos\phi=4\div5=0.8$、$P=100\times20\times0.8=\qty{1600}{W}$。抵抗の熱は $RI^2=4\times20^2=\qty{1600}{W}$。
\phtashikame
2つが同じ（コイルは平均して電力を使わない）。無効電力 $Q=VI\sin\phi=100\times20\times0.6=\qty{1200}{var}$ は $XI^2=3\times400$ と同じ。
\kotae \qty{20}{A}、力率0.8、\qty{1600}{W}（$=RI^2$）
\end{phreidai}
\begin{jibun}
抵抗\qty{6}{\ohm}とコイル（リアクタンス\qty{8}{\ohm}）の直列に\qty{100}{V}をつなぎます。電流、力率、使われる電力を求めましょう。
\bkjans{$Z=\qty{10}{\ohm}$、$I=\qty{10}{A}$、力率0.6、$P=\qty{600}{W}$（$=6\times10^2$）}
\end{jibun}

\label{ph:電-11:ex3}
\begin{phreidai}[モーターの力率をよくするコンデンサー]
単相\qty{200}{V}・\qty{50}{Hz}につないだモーターが、有効電力\qty{2}{kW}、力率0.7で動いています（数は仮のものとします）。電流と無効電力を求めましょう。また、全体の力率を1にするには、何\,\si{\micro F}のコンデンサーを並列につなげばよいですか。
\phbamen
工場のモーターに、力率改善のコンデンサーを選ぶ場面です。
\phmushi
コンデンサーの損失、電圧の波の形のくずれを無視します。モーターの動きは、コンデンサーをつないでも変わらないとします。
\phzu
$P=\qty{2000}{W}$、$\cos\phi=0.7$、$V=\qty{200}{V}$、$\omega=2\pi\times50\approx\qty{314}{rad/s}$。
\phriyu
電流は $P=VI\cos\phi$ から、無効電力は電力の三角形から、容量は⑤の $Q_C=\omega CV^2$ を $Q$ に等しくおいて求めます。
\phtateru
\[I=\frac{P}{V\cos\phi},\qquad Q=\sqrt{S^2-P^2}=P\tan\phi,\qquad C=\frac{Q}{\omega V^2}\]
\phtoku
\[I=\frac{2000}{200\times0.7}\approx\qty{14.3}{A},\qquad S=\frac{2000}{0.7}\approx\qty{2857}{VA}\]
\[Q=\sqrt{2857^2-2000^2}\approx\qty{2040}{var},\qquad C=\frac{2040}{314.2\times200^2}\approx\qty{1.62e-4}{F}\]
\phtashikame
単位：var÷(1/s×V$^2$)＝A·s/V＝C/V＝F。\\
大きさ：コンデンサーの電流は $\omega CV\approx\qty{10.2}{A}$。モーターの電流のおくれた分 $I\sin\phi\approx14.3\times0.714\approx\qty{10.2}{A}$ と同じ。
\kotae 約\qty{14}{A}、約\qty{2.0}{kvar}、約\qty{162}{\micro F}
\end{phreidai}
\begin{jibun}
単相\qty{100}{V}・\qty{60}{Hz}で、有効電力\qty{1.5}{kW}・力率0.75のモーターがあります。電流と無効電力、力率を1にするコンデンサーの容量を求めましょう。
\bkjans{\qty{20}{A}、約\qty{1.3}{kvar}（1323）、$C=1323\div(377\times100^2)$ で約\qty{351}{\micro F}}
\end{jibun}

\label{ph:電-11:ex4}
\begin{phreidai}[力率をよくすると、電線の損失が減る]
例題3で、コンデンサーをつないで力率を1にしました。電源からの電線を流れる電流は何 A になりますか。電線で熱になる損失は、もとの何倍になりますか。
\phbamen
力率改善の効き目を、電線の損失で確かめる場面です。
\phmushi
電線の抵抗は一定とし、電線での電圧の下がりは小さいとして無視します。
\phzu
有効電力 $P=\qty{2000}{W}$ は変わらない。$V=\qty{200}{V}$、改善後 $\cos\phi=1$。
\phriyu
電線の電流は $P/(V\cos\phi)$、損失は電流の2乗に比例します（⑥）。
\phtateru
\[I'=\frac{P}{V\times1},\qquad \frac{P_{\text{損}}'}{P_{\text{損}}}=\Bigl(\frac{I'}{I}\Bigr)^2\]
\phtoku
$I'=2000\div200=\qty{10}{A}$。損失の比は $(10\div14.29)^2=0.7^2=0.49$。
\phtashikame
比は力率の比の2乗 $(0.7/1)^2$ と同じ。モーターそのものの電流は\qty{14.3}{A}のまま（電線の電流だけが減る）。
\kotae \qty{10}{A}、損失は約0.49倍（約半分）
\end{phreidai}
\begin{jibun}
じぶんでとこう（例題3）のモーターで、力率を1にすると、電線の電流は何 A ですか。損失はもとの何倍ですか。
\bkjans{\qty{15}{A}、$(15\div20)^2$ で約0.56倍}
\end{jibun}

\label{ph:電-11:ex5}
\begin{phreidai}[高い電圧で送る]
発電所から\qty{100}{MW}の電力を、往復の抵抗が\qty{10}{\ohm}の送電線で送ります（数は仮のものとし、単相として考えます）。（1）力率1で、電圧\qty{100}{kV}と\qty{500}{kV}のとき、電流と損失（送る電力の何\,\%）を求めましょう。（2）\qty{100}{kV}で、力率が0.9なら、電流と損失は（1）の何倍ですか。
\phbamen
電-0 のはてな?（なぜ高い電圧で送る？）を、式で確かめる場面です。
\phmushi
送電線のコイルとコンデンサーの性質（電-13）、電圧の下がりを無視します。三相の3本を、単相の往復の抵抗にまとめて考えます。
\phzu
$P=\qty{1.0e8}{W}$、$R=\qty{10}{\ohm}$、$V=\qty{1.0e5}{V}$・$\qty{5.0e5}{V}$。
\phriyu
⑥の $I=P/(V\cos\phi)$ と $P_{\text{損}}=RI^2$ を使います。
\phtateru
\[I=\frac{P}{V\cos\phi},\qquad P_{\text{損}}=RI^2\]
\phtoku
（1）\qty{100}{kV}：$I=10^8\div10^5=\qty{1000}{A}$、$P_{\text{損}}=10\times1000^2=\qty{1.0e7}{W}$（10\,\%）。\\
\phantom{（1）}\qty{500}{kV}：$I=\qty{200}{A}$、$P_{\text{損}}=10\times200^2=\qty{4.0e5}{W}$（0.4\,\%）。\\
（2）電流は $1/0.9\approx1.11$ 倍、損失は $1/0.9^2\approx1.23$ 倍。
\phtashikame
電圧5倍で損失は $1/25$（$10\,\%\div25=0.4\,\%$）。力率が下がると、同じ電力に余分な電流が要り、損失は力率の2乗に反比例して増える。
\kotae （1）\qty{100}{kV}：\qty{1000}{A}・\qty{10}{MW}（10\,\%）、\qty{500}{kV}：\qty{200}{A}・\qty{0.4}{MW}（0.4\,\%）。（2）電流約1.1倍、損失約1.2倍
\end{phreidai}
\begin{jibun}
例題5の送電線で、\qty{275}{kV}・力率1で送ると、電流と損失（送る電力の何\,\%）はいくらですか。
\bkjans{電流約\qty{364}{A}、損失約\qty{1.3}{MW}（約1.3\,\%）}
\end{jibun}

\label{ph:電-11:ex6}
\begin{phreidai}[変圧器の大きさは kVA]
\qty{50}{kVA}の変圧器に、力率0.8の負荷をつなぎます。有効電力は何 kW まで使えますか。力率を1にすると、何 kW まで使えますか。
\phbamen
工場で、変圧器の大きさが足りるかを調べる場面です。
\phmushi
変圧器の中の損失、電圧の変化を無視します。
\phzu
皮相電力の上限 $S=\qty{50}{kVA}$、$\cos\phi=0.8$。
\phriyu
変圧器の電線の温まり方は電流（$S=VI$）で決まるので、上限は皮相電力で決まります（④）。
\phtateru
\[P=S\cos\phi\]
\phtoku
$P=50\times0.8=\qty{40}{kW}$。力率1なら $P=50\times1=\qty{50}{kW}$。
\phtashikame
力率をよくすれば、同じ変圧器でより多くの機械を動かせる（変圧器を大きくしなくてすむ）。
\kotae \qty{40}{kW}まで（力率1なら\qty{50}{kW}まで）
\end{phreidai}
\begin{jibun}
\qty{30}{kVA}の変圧器に、力率0.85の負荷をつなぐと、有効電力は何 kW まで使えますか。
\bkjans{\qty{25.5}{kW}}
\end{jibun}

\label{ph:電-11:ex7}
\begin{phreidai}[単相3線式の中性線]
家の単相3線式（両はしの線と中性線の間がそれぞれ\qty{100}{V}）で、片側（A 側）に\qty{1500}{W}の電子レンジ、もう片側（B 側）に\qty{1200}{W}のドライヤーをつなぎます（どちらも力率1とします）。A 側の線、B 側の線、中性線を流れる電流を求めましょう。両側が\qty{1000}{W}ずつなら、中性線の電流はいくらですか。
\phbamen
家の分電盤から来る3本の線の電流を調べる場面です。
\phmushi
電線の抵抗、電圧の下がりを無視します。
\phzu
A 側の線から中性線へ、中性線から B 側の線へ、という向きに同じ時刻の電流を考えます。A 側と B 側の電圧は、中性線から見て反対向き（180°ずれ）です。
\phriyu
中性線の点で、電流の和が0（電流則：電-4 ⑦）。A 側の電流は中性線へ入り、B 側の電流は中性線から出ていくので、中性線には差だけが流れます。
\phtateru
\[I_{\mathrm A}=\frac{P_{\mathrm A}}{100},\qquad I_{\mathrm B}=\frac{P_{\mathrm B}}{100},\qquad I_{\mathrm N}=|I_{\mathrm A}-I_{\mathrm B}|\]
\phtoku
$I_{\mathrm A}=\qty{15}{A}$、$I_{\mathrm B}=\qty{12}{A}$、$I_{\mathrm N}=\qty{3}{A}$。両側が\qty{1000}{W}ずつなら $10-10=0$。
\phtashikame
両側がつり合えば、中性線は要らない（三相の中性線と同じ考え）。
\kotae A 側\qty{15}{A}、B 側\qty{12}{A}、中性線\qty{3}{A}。つり合えば0
\end{phreidai}
\begin{jibun}
例題7で、A 側に\qty{1200}{W}、B 側に\qty{300}{W}をつなぐと、3本の線の電流はいくらですか。
\bkjans{A 側\qty{12}{A}、B 側\qty{3}{A}、中性線\qty{9}{A}}
\end{jibun}

\begin{phmijika}[中性線が切れると、あぶない]
例題7のような家で、もし中性線が外れると、\qty{1500}{W}の器具と\qty{500}{W}の器具が、\qty{200}{V}に直列につながった形になります。抵抗は $100^2\div1500\approx\qty{6.7}{\ohm}$ と $100^2\div500=\qty{20}{\ohm}$ なので、電圧は抵抗に比例して\qty{50}{V}と\qty{150}{V}に分かれ、小さい器具に\qty{150}{V}がかかって、こわれたり燃えたりするおそれがあります。分電盤には、これを見つけて電気を止める装置がついていることがあります。配線の工事は、資格をもつ人が行います。
\end{phmijika}

\label{ph:電-11:ex8}
\begin{phreidai}[三相\qty{200}{V}の1相の電圧]
工場の三相\qty{200}{V}（線間電圧）を、Y 結線の負荷につなぎます。1相の負荷にかかる電圧は何 V ですか。
\phbamen
三相の機械の、1つのコイルにかかる電圧を求める場面です。
\phmushi
電線の抵抗、3つの相の電圧の少しの違いを無視します（平衡とする）。
\phzu
$V=\qty{200}{V}$。求めるのは相電圧 $V_{\mathrm p}$。
\phriyu
Y 結線では、線間電圧は相電圧の $\sqrt3$ 倍です（⑦）。
\phtateru
\[V_{\mathrm p}=\frac{V}{\sqrt3}\]
\phtoku
$V_{\mathrm p}=200\div1.732\approx\qty{115}{V}$。
\phtashikame
$V_{\mathrm p}$ は $V$ より小さく、$V/2=100$ より大きい（120°ずれなので、差は2倍より小さい）。
\kotae 約\qty{115}{V}
\end{phreidai}
\begin{jibun}
線間電圧\qty{400}{V}の三相を Y 結線の負荷につなぐと、相電圧は何 V ですか。
\bkjans{約\qty{231}{V}}
\end{jibun}

\label{ph:電-11:ex9}
\begin{phreidai}[$\Delta$ の負荷を Y に置きかえる]
1相\qty{30}{\ohm}の抵抗3つを $\Delta$ 結線にして、線間電圧\qty{200}{V}の三相につなぎます。（1）相電流と線電流を求めましょう。（2）同じ線電流になる Y 結線の1相の抵抗を求めましょう。
\phbamen
ヒーターのつなぎ方を、計算しやすい Y の形に置きかえる場面です。
\phmushi
電線の抵抗を無視します。3つの抵抗は等しい（平衡）とします。
\phzu
$Z_\Delta=\qty{30}{\ohm}$、$V=\qty{200}{V}$。
\phriyu
$\Delta$ では1相に線間電圧がかかり、線電流は相電流の $\sqrt3$ 倍。Y では1相に $V/\sqrt3$ がかかり、線電流＝相電流です（⑦）。
\phtateru
\[I_{\mathrm p}=\frac{V}{Z_\Delta},\quad I=\sqrt3I_{\mathrm p}\qquad\text{と}\qquad I=\frac{V/\sqrt3}{Z_{\mathrm Y}}\]
\phtoku
（1）$I_{\mathrm p}=200\div30\approx\qty{6.67}{A}$、$I=1.732\times6.67\approx\qty{11.5}{A}$。\\
（2）2つの $I$ を等しくおくと $\sqrt3V/Z_\Delta=V/(\sqrt3Z_{\mathrm Y})$、$Z_{\mathrm Y}=Z_\Delta/3=\qty{10}{\ohm}$。
\phtashikame
Y の\qty{10}{\ohm}では、$115.5\div10\approx\qty{11.5}{A}$ で同じ。電力も $3\times6.67^2\times30\approx\qty{4.0}{kW}$、$3\times11.5^2\times10\approx\qty{4.0}{kW}$ で同じ。
\kotae （1）相電流 約\qty{6.7}{A}、線電流 約\qty{11.5}{A}。（2）\qty{10}{\ohm}
\end{phreidai}
\begin{jibun}
1相\qty{60}{\ohm}の $\Delta$ 結線の負荷を、線間電圧\qty{200}{V}の三相につなぎます。同じはたらきの Y の1相の抵抗と、線電流を求めましょう。
\bkjans{\qty{20}{\ohm}、線電流 約\qty{5.8}{A}}
\end{jibun}

\label{ph:電-11:ex10}
\begin{phreidai}[三相のモーターの電力]
線間電圧\qty{200}{V}の三相のモーターに、線電流\qty{10}{A}が流れ、力率は0.85です。モーターが受けとる電力を求めましょう。
\phbamen
工場の三相のモーターの電力を、線ではかった電圧・電流から求める場面です。
\phmushi
3つの相のわずかな違いを無視し、平衡とします。
\phzu
$V=\qty{200}{V}$、$I=\qty{10}{A}$、$\cos\phi=0.85$。
\phriyu
線間電圧と線電流で書いた三相の電力の式（⑧）を使います。Y か $\Delta$ かを知らなくても使えます。
\phtateru
\[P=\sqrt3VI\cos\phi\]
\phtoku
$P=1.732\times200\times10\times0.85\approx\qty{2944}{W}$。
\phtashikame
$3VI\cos\phi$（約\qty{5.1}{kW}）としてしまうと、$\sqrt3$ 倍の大きすぎる答えになる（まちがえやすいところ）。Y なら1相 $115.5\times10\times0.85=\qty{981}{W}$ の3倍で同じ。
\kotae 約\qty{2.9}{kW}
\end{phreidai}
\begin{jibun}
線間電圧\qty{400}{V}、線電流\qty{20}{A}、力率0.9の三相の負荷の電力を求めましょう。
\bkjans{約\qty{12.5}{kW}}
\end{jibun}

\label{ph:電-11:ex11}
\begin{phreidai}[三相の瞬時電力を Python で]
相電圧\qty{100}{V}・相電流\qty{10}{A}（実効値）・力率0.85の平衡三相の負荷で、3つの相の瞬時電力の和を、$\omega t$ を60°ごとに計算しましょう。
\phbamen
⑧の「三相の電力は一定」を、数で確かめる場面です。
\phmushi
電圧・電流は正弦波で、3つの相はちょうど120°ずつずれているとします。
\phzu
$v_k=\sqrt2\cdot100\sin(\omega t-120°k)$、$i_k=\sqrt2\cdot10\sin(\omega t-120°k-\phi)$（$k=0,1,2$）、$\cos\phi=0.85$。
\phriyu
瞬時電力は $p=\sum v_ki_k$。各時刻で数を入れて計算すれば、式を変形しなくても確かめられます。
\phtateru
\begin{phcode}[Python]
import math
V, I, pf = 100, 10, 0.85
phi = math.acos(pf)
for deg in range(0, 361, 60):
    t = math.radians(deg)
    p = 0
    for k in range(3):
        s = k*2*math.pi/3
        v = math.sqrt(2)*V*math.sin(t - s)
        i = math.sqrt(2)*I*math.sin(t - s - phi)
        p = p + v*i
    print(deg, round(p, 1))
\end{phcode}
\phtoku
\begin{phcodeout}
0 2550.0
60 2550.0
120 2550.0
180 2550.0
240 2550.0
300 2550.0
360 2550.0
\end{phcodeout}
どの時刻も\qty{2550}{W}です。
\phtashikame
$3V_{\mathrm p}I_{\mathrm p}\cos\phi=3\times100\times10\times0.85=2550$ と同じ。1相だけ（\texttt{range(1)}）にすると、$-150$〜\qty{1850}{W}の間でゆれる。
\kotae どの時刻も\qty{2550}{W}で一定
\end{phreidai}
\begin{jibun}
例題11のプログラムで、力率を0.5にすると、和はいくらになりますか。1相だけの瞬時電力は、いくらからいくらの間でゆれますか。
\bkjans{いつも\qty{1500}{W}。1相は $500-1000=\qty{-500}{W}$ から $500+1000=\qty{1500}{W}$ の間}
\end{jibun}

\label{ph:電-11:ex12}
\begin{phreidai}[三相は銅が少なくてすむ（はてな？）]
同じ電力 $P$ を、同じ線間電圧 $V$・力率1で、同じ損失になるように送ります。単相2線式と三相3線式で、要る電線の銅の量（体積）を比べましょう。電線の長さは同じとします。
\phbamen
送電に三相を使う理由を、電線の材料の量で確かめる場面です。
\phmushi
電線の抵抗だけを考え、コイルの性質などは無視します。銅の量は、1本の断面積×本数に比例し、1本の抵抗は断面積に反比例するとします（$R=\rho l/S$：電-4）。
\phzu
単相：2本、1本の抵抗 $R_1$。三相：3本、1本の抵抗 $R_3$。
\phriyu
損失を等しくおいて $R_3$ と $R_1$ の関係を出し、銅の量（本数÷1本の抵抗に比例）を比べます。
\phtateru
\[\text{単相：}I_1=\frac{P}{V},\ P_{\text{損}}=2R_1I_1^2\qquad\text{三相：}I_3=\frac{P}{\sqrt3V},\ P_{\text{損}}=3R_3I_3^2\]
\phtoku
三相の損失は $3R_3\times\dfrac{P^2}{3V^2}=R_3\dfrac{P^2}{V^2}$、単相は $2R_1\dfrac{P^2}{V^2}$。等しいので $R_3=2R_1$。銅の量の比は
\[\frac{3/R_3}{2/R_1}=\frac{3}{2\times2}=\frac34\]
\phtashikame
三相は線が1.5倍の本数だが、1本の電流が $1/\sqrt3$ なので、1本を細く（抵抗2倍）できる。Python で数を入れても0.75。
\kotae 三相は単相の $\tfrac34$ の銅ですむ
\end{phreidai}
\begin{jibun}
銅の量を同じにして、単相2線式と三相3線式で同じ電力を同じ線間電圧で送ると、三相の損失は単相の何倍ですか。
\bkjans{三相の1本は $R_3=1.5R_1$。損失の比は $R_3/(2R_1)=\tfrac34$ 倍}
\end{jibun}

\label{ph:電-11:ex13}
\begin{phreidai}[回転磁界の速さ]
\qty{50}{Hz}の三相交流で、2極と4極の回転磁界をつくります。それぞれ、1分に何回転しますか。
\phbamen
工場のモーターの回る速さの見当をつける場面です（電-12 の誘導モーター・同期モーター）。
\phmushi
コイルの形による磁場のむらを無視し、磁場はなめらかに回るとします。
\phzu
$f=\qty{50}{Hz}$、極の数2と4。
\phriyu
2極の回転磁界は1周期で1回転、4極は半回転です（⑨）。
\phtateru
\[n_{\mathrm s}=\frac{120f}{\text{極の数}}\]
\phtoku
2極：$120\times50\div2=\qty{3000}{min^{-1}}$。4極：$120\times50\div4=\qty{1500}{min^{-1}}$。
\phtashikame
2極は1秒に50回転（$50\times60=3000$）。極の数を2倍にすると、速さは半分。
\kotae 2極は毎分3000回転、4極は毎分1500回転
\end{phreidai}
\begin{jibun}
\qty{60}{Hz}の三相交流で、2極と6極の回転磁界の速さを求めましょう。
\bkjans{2極：毎分3600回転、6極：毎分1200回転}
\end{jibun}

%% ======================== まちがえやすいところ ========================
\Needspace{11\baselineskip}\bksection{tsumazuki}{まちがえやすいところをチェック}
\begin{tsumazuki}{交流の電力は「電圧×電流」？}
\bkng{\qty{100}{V}で\qty{10}{A}流れるので、使われる電力は\qty{1000}{W}。}
\bkok{力率がわからないと決まらない。力率0.8なら\qty{800}{W}。\qty{1000}{VA}は皮相電力。}
わけ：電流のうち電圧と山がずれた分は、エネルギーを行き来させるだけで、使われないからです。
\end{tsumazuki}

\begin{tsumazuki}{3つの電力をたす}
\bkng{有効電力\qty{800}{W}と無効電力\qty{600}{var}をたして、皮相電力は\qty{1400}{VA}。}
\bkok{電力の三角形の辺なので $\sqrt{800^2+600^2}=\qty{1000}{VA}$。}
わけ：$P$ と $Q$ は、90°ずれた向きの量です（フェーザの直角な2つの向き）。
\end{tsumazuki}

\begin{tsumazuki}{三相の電力}
\bkng{線間電圧\qty{200}{V}・線電流\qty{10}{A}・力率1の三相の電力は、3倍して $3\times200\times10=\qty{6000}{W}$。}
\bkok{$\sqrt3\times200\times10\approx\qty{3460}{W}$。3倍してよいのは、相電圧×相電流のとき。}
わけ：Y 結線なら、1相の電圧は $200/\sqrt3$ です。
\end{tsumazuki}

\begin{tsumazuki}{コンデンサーで力率改善すると}
\bkng{コンデンサーをつなぐと、モーターの電流が減って、モーターの力も弱くなる。}
\bkok{モーターの電流も仕事も変わらない。減るのは、電源からの電線を流れる電流。}
わけ：コンデンサーとモーターの間で無効電力をやりとりし、電源まで行き来する分がなくなるからです。
\end{tsumazuki}

\begin{tsumazuki}{電流のおくれと、電力の符号}
\bkng{コイルの電力の図で、$p$ がマイナスの時刻は、コイルが電気を使わずに捨てている。}
\bkok{マイナスは、コイルがためた磁場のエネルギーを、電源へ返している時刻。}
わけ：$p=vi$ の符号は、エネルギーの流れの向きを表します。
\end{tsumazuki}

%% ======================== 工学ではこう使う ========================
\Needspace{18\baselineskip}\bksection{phgear}{工学ではこう使う}
\begin{phkougaku}[電力をはかる・料金を決める・力率を保つ]
\textbf{電力量計（電気のメーター）}：家や工場のメーターが数えるのは、有効電力を時間でたした電力量（kWh）です。無効電力の行き来は、平均すると0なので数えません。電子式のメーターは、電圧と電流を短い時間おきにはかり、その積 $vi$ を足しあわせて電力量を求めます（③の平均を、計算機で行っている）。

\textbf{力率と電気料金}：工場など高い電圧で電気を受ける契約では、力率が0.85を上回ると、上回る1\,\%ごとに基本料金を1\,\%割り引く、というしくみの料金表があります。無効電力が少ないと、電力会社の電線や変圧器が楽になるからです。工場は、受電の設備に進相コンデンサーを置き、力率を1に近く保ちます。

\textbf{三相のモーターと発電機}：発電所の発電機は三相で、つくられた電力は回転磁界のしくみでそのまま三相のモーターを回せます（電-12）。電池の直流や、単相の交流からでも、インバーター（電-12 ⑨）で三相の交流をつくれば、三相のモーターを回せます。

\begin{phketa}[電圧と電力の数の目安]
\phketaline{家のコンセント}{\qty{100}{V}は\qtyrange{95}{107}{V}、\qty{200}{V}は\qtyrange{182}{222}{V}に保つ（法令）}
\phketaline{柱上変圧器}{\qty{6600}{V}を\qty{210}{V}・\qty{105}{V}に下げる}
\phketaline{配電・送電}{配電\qty{6.6}{kV}、送電\qtyrange{66}{500}{kV}}
\phketaline{送電線などの磁場の基準}{人のいる所の平均\qty{200}{\micro T}以下（法令、商用周波数）}
\phboxhead{答えが桁で合っているか}
家の電気製品の電流は数 A〜十数 A（契約は数十 A）。工場の三相のモーターは数 kW〜数百 kW。送電線の電流は数百〜数千 A。損失が送る電力の何十\,\% にもなる答えが出たら、電圧の単位（kV と V）を取り違えていないか確かめましょう。
\end{phketa}
\end{phkougaku}
\phstation{電力工学,電気設備,電気機器,電力系統工学}{有効電力、無効電力、皮相電力、力率、力率改善、進相コンデンサー、三相交流、線間電圧、相電圧、Y 結線・$\Delta$ 結線、回転磁界、対称座標法}

%% ======================== 章末問題 ========================
\bkendproblems
\bkproblemgroup{基本}{3}
\bkproblem \qty{200}{V}の交流で、\qty{15}{A}・力率0.8の機械の有効電力・無効電力・皮相電力を求めましょう。
\bkpans{$P=\qty{2.4}{kW}$、$Q=\qty{1.8}{kvar}$、$S=\qty{3.0}{kVA}$}
\bkproblem 抵抗\qty{8}{\ohm}とコイル（リアクタンス\qty{6}{\ohm}）の直列に\qty{100}{V}の交流をつなぎます。電流・力率・有効電力・無効電力を求めましょう。
\bkpans{$Z=\qty{10}{\ohm}$、$I=\qty{10}{A}$、力率0.8、$P=\qty{800}{W}$（$=8\times10^2$）、$Q=\qty{600}{var}$}
\bkproblem 単相\qty{200}{V}・\qty{50}{Hz}で、\qty{4}{kW}・力率0.8の負荷があります。（1）電流を求めましょう。（2）力率を1にするコンデンサーの容量を求めましょう。（3）そのとき電線の電流は何 A になり、損失はもとの何倍になりますか。
\bkpans{（1）\qty{25}{A}。（2）$Q=\qty{3}{kvar}$、$C=3000\div(314.2\times200^2)$ で約\qty{239}{\micro F}。（3）\qty{20}{A}、0.64倍}
\bkproblem 同じ電力を、同じ送電線・力率1で送るとき、電圧を\qty{66}{kV}から\qty{154}{kV}に上げると、損失はもとの何倍になりますか。
\bkpans{$(66\div154)^2$ で約0.18倍（約5分の1）}
\bkproblem \qty{100}{V}・\qty{5}{A}・力率0.6の機械を2時間使いました。電力量計が数える電力量は何 kWh ですか。
\bkpans{$100\times5\times0.6\times2=\qty{600}{Wh}$ で、\qty{0.6}{kWh}}
\bkproblem 単相3線式の家で、A 側（\qty{100}{V}）に\qty{1000}{W}、B 側（\qty{100}{V}）に\qty{600}{W}、両はしの線の間（\qty{200}{V}）に\qty{2000}{W}の器具をつなぎます（どれも力率1）。A 側の線・B 側の線・中性線の電流を求めましょう。
\bkpans{\qty{200}{V}の器具は\qty{10}{A}で、両はしの線だけを流れる。A 側の線 $10+10=\qty{20}{A}$、B 側の線 $6+10=\qty{16}{A}$、中性線 $10-6=\qty{4}{A}$}
\bkproblemgroup{チャレンジ}{5}
\bkproblem 1つ\qty{20}{\ohm}のヒーター3つを、線間電圧\qty{200}{V}の三相につなぎます。（1）$\Delta$ 結線のときの相電流・線電流・電力、（2）Y 結線のときの線電流・電力を求めましょう。つなぎ方を変えると、電力は何倍になりますか。
\bkpans{（1）相電流\qty{10}{A}、線電流 約\qty{17.3}{A}、\qty{6.0}{kW}。（2）線電流 約\qty{5.8}{A}、\qty{2.0}{kW}。Y は $\Delta$ の $\tfrac13$}
\bkproblem 線間電圧\qty{200}{V}の三相で、\qty{30}{kW}・力率0.75の工場があります。（1）線電流を求めましょう。（2）力率を0.95にするには、コンデンサーで何 kvar の無効電力を打ち消せばよいですか。（3）そのとき線電流は何 A になり、電線の損失は何倍になりますか。
\bkpans{（1）$30000\div(\sqrt3\times200\times0.75)$ で約\qty{115}{A}。（2）$Q$ は $30\times\tan(\cos^{-1}0.75)\approx\qty{26.5}{kvar}$ から $30\times\tan(\cos^{-1}0.95)\approx\qty{9.9}{kvar}$ へ。約\qty{16.6}{kvar}。（3）約\qty{91}{A}、$(0.75\div0.95)^2$ で約0.62倍}
\bkproblem ⑨の3つのコイルで、1つのコイルが電流の山のときにつくる磁場を $B_0=\qty{0.010}{T}$ とします。$\omega t=30°$ のときの3つのコイルの磁場（大きさと向き）を求め、和が $\tfrac32B_0$ で30°の向きになることを、成分で確かめましょう。
\bkpans{a：$0.010\cos30°\approx\qty{0.0087}{T}$（0°の向き）。b：$\cos(-90°)=0$。c：$0.010\cos(-210°)\approx-0.0087$ で、240°の向きの反対の60°の向きに\qty{0.0087}{T}。和の横は $0.0087+0.0087\cos60°\approx0.0130$、たては $0.0087\sin60°=0.0075$。大きさ $\sqrt{0.0130^2+0.0075^2}\approx\qty{0.015}{T}$、向きは $\tan^{-1}(0.0075\div0.0130)=30°$}
\bkproblem 中性線のある三相（Y 結線）で、3つの相の抵抗の負荷に、a 相\qty{10}{A}、b 相\qty{10}{A}、c 相\qty{4}{A}の電流が流れています（どれも力率1、相の電圧は120°ずつずれる）。中性線の電流を、フェーザ（複素数）のたし算で求めましょう。
\bkpans{$10+10\bigl(-\tfrac12-\tfrac{\sqrt3}{2}j\bigr)+4\bigl(-\tfrac12+\tfrac{\sqrt3}{2}j\bigr)=3-3\sqrt3j$。大きさ $\sqrt{9+27}=\qty{6}{A}$}

%% ======================== ふりかえり ========================
\begin{furikaeri}
\begin{checklist}
\item 交流の電力が時間とともにゆれること、家の電気と送電線の電気の違いを、図で説明できる。\bkfurref{ph:電-11:sec2}
\item 力率を使って、有効電力・無効電力・皮相電力を計算できる。\bkfurref{ph:電-11:sec4}
\item コンデンサーで力率をよくすると、電線の電流と損失が減るわけを説明できる。\bkfurref{ph:電-11:sec5}
\item 三相交流の Y 結線・$\Delta$ 結線の電圧と電流、三相の電力を計算できる。\bkfurref{ph:電-11:sec7}
\item 三相の電力が時間によらず一定になること、回転磁界ができることを、式で示せる。\bkfurref{ph:電-11:sec9}
\end{checklist}
\end{furikaeri}
\bkfukidashi{mina}{3本の電流の和がいつも0だから、帰りの線が要らないんだね。ぐるぐる回る磁場まで作れるなんて、120°ずらすのはすごい工夫だ！}
